% paper4_unification.tex
% Temporal Asymmetry as Organizing Principle:
% From Quantum Channels to Cosmological Structure
% Paper IV in the tau framework series
% Author: Sheng-Kai Huang
% Compile with: pdflatex paper4_unification.tex

\documentclass[prd,twocolumn,superscriptaddress,longbibliography,floatfix]{revtex4-2}

\usepackage{amsmath}
\usepackage{amssymb}
\usepackage{amsfonts}
\usepackage{amsthm}
\usepackage{bm}
\usepackage{hyperref}
\usepackage{booktabs}
\usepackage{graphicx}

\newtheorem{theorem}{Theorem}
\newtheorem{corollary}[theorem]{Corollary}
\newtheorem{conjecture}[theorem]{Conjecture}
\newtheorem{postulate}{Postulate}
\newtheorem*{proposition}{Proposition}

% Convenience macros (consistent with Papers 1--3, 5)
\newcommand{\kB}{k_{\mathrm{B}}}
\newcommand{\Tr}{\mathrm{Tr}}
\newcommand{\id}{\mathrm{id}}
\newcommand{\Petz}{\mathcal{R}}
\newcommand{\chan}{\mathcal{N}}
\newcommand{\hilb}{\mathcal{H}}
\newcommand{\code}{\mathcal{C}}
\newcommand{\KL}{D}
\newcommand{\ket}[1]{|#1\rangle}
\newcommand{\bra}[1]{\langle#1|}
\newcommand{\braket}[2]{\langle#1|#2\rangle}
\newcommand{\op}[2]{|#1\rangle\langle#2|}
\newcommand{\abs}[1]{\left|#1\right|}
\newcommand{\norm}[1]{\left\|#1\right\|}
\newcommand{\tauasym}{\tau}

% Paper 2 macros
\newcommand{\ceff}{c_{\mathrm{eff}}}
\newcommand{\rs}{r_{\mathrm{s}}}
\newcommand{\Sgrav}{\Sigma_{\mathrm{grav}}}
\newcommand{\Fbound}{F_{\mathrm{bound}}}
\newcommand{\nref}{n_{\mathrm{grav}}}

% Paper 4 specific macros
\newcommand{\Daraki}{D_{\mathrm{Araki}}}
\newcommand{\Aobs}{\mathcal{A}_{\mathcal{O}}}
\newcommand{\Atot}{\mathcal{A}_{\mathrm{tot}}}
\newcommand{\Ast}{\mathcal{A}_{\mathrm{st}}}
\newcommand{\Am}{\mathcal{A}_{\mathrm{m}}}
\newcommand{\rhost}{\rho_{\mathrm{st}}}
\newcommand{\rhom}{\rho_{\mathrm{m}}}
\newcommand{\SigO}{\Sigma_{\mathcal{O}}}
\newcommand{\tauO}{\tauasym_{\mathcal{O}}}
\newcommand{\Tmod}{T_{\mathrm{mod}}}
\newcommand{\GN}{G_{\mathrm{N}}}
\newcommand{\SdS}{S_{\mathrm{dS}}}
\newcommand{\rA}{r_{\mathrm{A}}}
\newcommand{\OmDM}{\Omega_{\mathrm{DM}}}

\begin{document}

\title{Temporal Asymmetry as Organizing Principle:\\
From Quantum Channels to Cosmological Structure}

\author{Sheng-Kai Huang}
\email{akai@fawstudio.com}
\affiliation{Independent Researcher}

\date{\today}

\begin{abstract}
The temporal asymmetry parameter $\tauasym = 1 - F$,
where $F$ is the Petz recovery fidelity, was shown in
Papers~I--III and V to govern quantum information recovery,
strong-field gravity, galactic dynamics, and observer-dependent
time, respectively.  Here we propose that these are not four
separate applications but five limits of a single equation:
\begin{equation*}
\Sigma \;=\; \Daraki\!\bigl(\omega\big|_{\Aobs}
\;\big\|\; \omega_{\mathrm{ref}}\big|_{\Aobs}\bigr)\,,
\end{equation*}
the Araki relative entropy restricted to the observer's
accessible algebra $\Aobs$.  The observer $\mathcal{O}$
selects the subalgebra, and different choices reproduce
(a)~channel QRE in flat spacetime [\textsc{proved}],
(b)~$\Sgrav = -\ln(-g_{00})$ for static gravity
[\textsc{derived}: gravitational thermal attenuator,
Paper~II],
(c)~running-$G$ rotation curves [\textsc{derived}: DPI
$+$ extensivity $\Rightarrow$ $\eta = 1$, Paper~III],
(d)~the Friedmann equations from entanglement equilibrium
[\textsc{proved}], and
(e)~observer-dependent $\tauO$ [\textsc{proved}].
We introduce a new postulate---Observer Entanglement
Equilibrium (OEE), $\delta\Sigma/\delta u^a = 0$---which promotes
the observer's 4-velocity to a dynamical field structurally
identical to the Khronon of Blanchet--Skordis
theory~\cite{BlanchetSkordis2024}.  The Petz-optimality
principle suggests a quadratic kinetic function $K(Q) \propto
(Q-1)^2$, yielding a non-propagating scalar mode ($c_s = 0$) that
mimics cold dark matter at the perturbative level.  We propose
two numerical relations that pin down the Khronon mass
$\mu$ in terms of cosmological observables:
$\mu c^2 = a_0(1+z_{\mathrm{dec}})$ and
$\mu = 2\pi(1+\Omega_{\mathrm{b}})/r_{\mathrm{d}}$,
each reproducing the Blanchet--Skordis fitted value
$\mu^{-1} = 22.3\,\mathrm{Mpc}$ to better than $0.3\%$.
Their mutual consistency yields a parameter-free prediction
for the MOND acceleration:
$a_0 = 2\pi(1+\Omega_{\mathrm{b}})c^2/
[r_{\mathrm{d}}(1+z_{\mathrm{dec}})]
= 1.197 \times 10^{-10}\,\mathrm{m\,s}^{-2}$
($0.3\%$ match), a $30$--$40\times$ improvement over
Milgrom's $a_0 = cH_0/(2\pi)$.
The old identification $\mu_0 = H_0/c$, proposed in
an earlier version, is excluded by CMB data at $>55\sigma$.
The dark matter fraction $\OmDM \sim \rho_{\mathrm{crit}}/3$
remains an order-of-magnitude estimate ($\sim$25\% above
observed), with the exact $O(1)$ prefactor undetermined.
We clearly delineate what is proved, conjectured, and
speculative, identify six falsifiable predictions unique to the
unified framework, and list the open problems whose resolution
would promote this organizational principle to a predictive theory.
\end{abstract}

\maketitle

%=======================================================================
\section{Introduction}
\label{sec:intro}
%=======================================================================

The temporal asymmetry parameter
\begin{equation}
\tauasym \;\equiv\; 1 - F\!\bigl(\rho,\;
\Petz_{\sigma,\chan}\!\circ\!\chan(\rho)\bigr)
\;\in\; [0,1]
\label{eq:tau_def}
\end{equation}
was introduced in Paper~I~\cite{Paper1} as a unified measure
of irreversibility, bounded by
\begin{equation}
F \;\geq\; e^{-\Sigma/2}\,,
\label{eq:JRSWW_bound}
\end{equation}
where $\Sigma = \KL(\rho\|\sigma) -
\KL\!\bigl(\chan(\rho)\|\chan(\sigma)\bigr)$ is the entropy
production of the channel $\chan$, and $\Petz_{\sigma,\chan}$
is the Petz recovery map---the unique Bayesian retrodiction
functor~\cite{Parzygnat2023,Petz1988}.
Paper~I proved a four-way equivalence chain:
$\tauasym = 0$ if and only if $\Sigma = 0$, if and only if the
state forms a quantum Markov chain, if and only if the
Knill--Laflamme conditions for quantum error correction are
satisfied~\cite{Paper1}.

Paper~II~\cite{Paper2} identified the gravitational entropy
production as $\Sgrav = -\ln(-g_{00})$ and established the
triple identification $\ceff/c = e^{-\Sgrav/2} = \Fbound$:
the coordinate speed of light \emph{is} the Petz recovery
fidelity bound.  Paper~III~\cite{Paper3} showed that
quantum corrections to the gravitational channel---the
renormalization-group running of Newton's
constant~\cite{Kumar2025}---produce logarithmic corrections
to $\Sigma$ that yield flat rotation curves without dark matter
particles.  Paper~V~\cite{Paper5} formalized the observation
that $\tauasym$ is observer-dependent:
different observers with access to different subalgebras
assign different temporal asymmetries to the same
process~\cite{BassoCeleri2025,DeVuyst2025}.

A natural question arises: are these four applications
independent, or are they limits of a single structure?
In this paper we argue that they are the latter.
We propose (as a new organizational framework) the
\emph{master equation}
\begin{equation}
\boxed{\;
\Sigma_{\mathcal{O}}
\;=\; \Daraki\!\bigl(\omega\big|_{\Aobs}
\;\big\|\; \omega_{\mathrm{ref}}\big|_{\Aobs}\bigr)
\;\;}
\label{eq:master}
\end{equation}
where $\omega$ is the global quantum state of the universe,
$\omega_{\mathrm{ref}}$ is a reference state (vacuum, Gibbs,
or matter-induced geometry), $\Aobs$ is the von Neumann
subalgebra accessible to observer $\mathcal{O}$, and
$\Daraki$ is the Araki relative entropy~\cite{Araki1976}.
The claim is not that this equation is new---each ingredient
is known---but that its different \emph{limits}, selected by
different observer algebras, reproduce the physics of
Papers~I--III and V.  This is analogous to how $F = ma$
yields Kepler orbits, simple harmonic motion, and
Navier--Stokes flow depending on the force law and
boundary conditions.

We also introduce a new postulate---\emph{Observer
Entanglement Equilibrium} (OEE)---which promotes
the observer's 4-velocity $u^a$ from an external
parameter to a dynamical field.  This produces a field
structurally identical to the Khronon of
Blanchet--Skordis theory~\cite{BlanchetSkordis2024},
providing a physical interpretation of the aether/Khronon
as ``the observer's time direction made dynamical.''

Throughout, we use a three-tier classification:
\textsc{proved} (rigorous mathematical theorem),
\textsc{conjectured} (well-motivated, supported by
evidence, but no complete proof), and
\textsc{speculative} (interesting direction, minimal
mathematical support).  We do not disguise conjectures
as theorems.


%=======================================================================
\section{The Master Equation}
\label{sec:master}
%=======================================================================

\subsection{Algebraic setting}

Let $\Atot$ be the von Neumann algebra of all observables
of the universe, generated by a spacetime subalgebra
$\Ast$ and a matter subalgebra $\Am$:
$\Atot = \Ast \vee \Am$.
A global state $\omega$ on $\Atot$ restricts to
$\rhost = \omega|_{\Ast}$ and $\rhom = \omega|_{\Am}$.
The \emph{observer} $\mathcal{O}$ is specified by a
subalgebra $\Aobs \subseteq \Atot$.  The restriction
$\omega|_{\Aobs}$ is the state accessible to
$\mathcal{O}$.

For Type~III algebras (generic in QFT), individual von
Neumann entropies diverge, but the Araki relative entropy
\begin{equation}
\Daraki(\omega\|\omega_0) \;=\;
-\bigl\langle\Omega,\;
\bigl(\ln\Delta_{\omega_0,\omega}\bigr)\Omega
\bigr\rangle
\label{eq:araki}
\end{equation}
(where $\Delta_{\omega_0,\omega}$ is the relative modular
operator and $\Omega$ represents $\omega$ in the natural
cone) is well-defined, non-negative, and satisfies the
data-processing inequality (DPI)~\cite{Araki1976,Lindblad1975}.
This is why $\Sigma$---not entropy $S$---is the
fundamental quantity.

When the observer carries a clock (quantum reference
frame), the Type~III algebra is promoted to Type~II via
a crossed-product construction, and the generalized
entropy becomes well-defined~\cite{Chandrasekaran2023,
DeVuyst2025}.  Different clocks yield different
crossed products and different generalized entropies:
this is the algebraic origin of observer-dependent
$\tauO$.

\subsection{Connection to existing frameworks}

The master equation~\eqref{eq:master} sits at the
intersection of two recent results.

\emph{Dorau--Much}~\cite{DorauMuch2025} derive the
semiclassical Einstein equations from the Araki relative
entropy:
\begin{equation}
S^{\mathrm{rel}}(\omega_0\|\omega_\phi) \;=\;
-2\pi\!\int_{\mathcal{H}}\!
U\,\langle{:}T_{UU}{:}\rangle\,dU\,d\mathrm{vol}_S\,,
\label{eq:dorau_much}
\end{equation}
combined with the Bekenstein--Hawking relation
$S^{\mathrm{rel}} = \delta A/(4\GN)$
and the Raychaudhuri equation.  This is
Eq.~\eqref{eq:master} restricted to the near-horizon
regime.

\emph{Bianconi}~\cite{Bianconi2025} constructs the
``Gravity from Entropy'' (GfE) action as a generalized
matrix relative entropy between the spacetime metric
$\tilde{g}$ and a matter-induced metric $\tilde{G}$:
\begin{equation}
S_{\mathrm{GfE}} \;=\;
\int\!\sqrt{-g}\,
\Tr\bigl(\tilde{g}\,\ln\tilde{G}^{-1}\bigr)\,d^4x\,.
\label{eq:bianconi}
\end{equation}
In the low-coupling limit, this reduces to the
Einstein--Hilbert action coupled to Klein--Gordon
matter.  The identification of the gravitational action
with a relative-entropy-like functional is the
mathematical content of our master equation, realized
in a specific way.

We emphasize an important caveat.  Bianconi's generalized
matrix relative entropy operates on metric tensors (whose
eigenvalues are not normalized), while the Araki relative
entropy operates on states of a von Neumann algebra.
The exact mathematical relationship between these two
constructions is an open problem.
Equation~\eqref{eq:master} should therefore be understood
as a \emph{principle}---the gravitational dynamics are
governed by a relative entropy between spacetime and matter
degrees of freedom---not as a specific mathematical formula
with a unique realization.

\subsection{Three action-independent routes to $\Sgrav$}

The identity $\Sgrav = -\ln(-g_{00})$ is supported by
three independent derivations that do not depend on any
specific action principle~\cite{Paper2}:

\emph{Route~1: Modular flow.}---The Bisognano--Wichmann
theorem identifies the modular operator with the boost
generator on the Rindler wedge.  Extending the Dorau--Much
framework to a Schwarzschild background, the fractional
QRE loss is $(D_{\mathrm{in}} - D_{\mathrm{out}})/
D_{\mathrm{in}} = \rs/r$, exact for
Schwarzschild~\cite{DorauMuch2025}.

\emph{Route~2: Gravitational Landauer principle.}---An
erasure at radius $r$ in a gravitational field costs
$W(r) = \kB T(r)\ln 2 = \kB T_\infty \ln 2/
\sqrt{-g_{00}}$, where the Tolman temperature gradient
increases the cost.  The net entropy production is
$\Sigma = -\ln(-g_{00})$~\cite{Herrera2020}.

\emph{Route~3: Pure-loss bosonic channel.}---Modeling
gravitational redshift as a pure-loss bosonic channel
with transmissivity $\eta = -g_{00}$, the QRE decrease
is $\Sigma = -\ln(\eta) = -\ln(-g_{00})$.

All three routes share the mathematical core
$\Sgrav = -\ln(-g_{00})$.  The exponential metric
follows if the Petz bound $F = e^{-\Sigma/2}$ is
exactly saturated, giving
$g_{00} = -e^{-\rs/r}$~\cite{Paper2,Yilmaz1958}.

\subsection{The observer selects the limit}

The key structural insight is that the observer
$\mathcal{O}$ determines which subalgebra to restrict
to, and this restriction selects the effective physics.
Different observers form a lattice under refinement:
\begin{equation}
\bot \;\leq\; \mathcal{O}_{\mathrm{thermo}}
\;\leq\; \mathcal{O}_{\mathrm{lab}}
\;\leq\; \cdots \;\leq\; \mathcal{O}_{\mathrm{total}}
\;=\; \top\,,
\label{eq:lattice}
\end{equation}
where $\bot$ is the trivial observer (no access) and
$\top$ is the total observer ($\Aobs = \Atot$,
$\tauasym = 0$).  By Theorem~\ref{thm:monotone},
moving up the lattice decreases $\tauO$.

The \emph{same} physical state $\omega$, restricted to
different subalgebras, gives different $\Sigma$ values
and therefore different effective physics.  The QI
community studies the $\chan$-induced $\Sigma$ in flat
spacetime; the GR community studies
$\Sgrav = -\ln(-g_{00})$ at the asymptotic boundary;
the cosmology community measures $\Sigma$ at galactic
and Hubble scales; the philosophy community examines the
observer-dependence itself.  All study the same
Eq.~\eqref{eq:master}, restricted to different algebras.

\subsection{The zeroth statement: Wheeler--DeWitt}

The Wheeler--DeWitt equation
$\hat{H}|\Psi\rangle = 0$ implies no time
evolution at the fundamental level---the ``block universe.''
In our language:
\begin{equation}
\tauasym_{\mathrm{total}} \;=\; 0\,.
\label{eq:zeroth}
\end{equation}
The universe, as a closed system, has no arrow of time.
This is Level~0, from which everything follows.

For the total observer ($\Aobs = \Atot$):
$\chan_{\mathrm{total}} = \id$, so $\Sigma = 0$ and
$\tauasym = 0$.  For \emph{any} partial observer
$\mathcal{O}$ with
$\Aobs \subsetneq \Atot$:
$\chan_{\mathcal{O}} = \Tr_{\Atot \setminus \Aobs}$,
and the DPI guarantees $\SigO \geq 0$.
The arrow of time is the \emph{cost} of being a
partial observer.

This connects to the Page--Wootters
mechanism~\cite{PageWootters1983}: conditioning on a
clock $C$ extracts effective time evolution from a
timeless state.  In our framework, the clock is one
observer choice $\mathcal{O}_C$, the conditioned
state is $\omega|_{\Aobs[C]}$, and the entropy
production $\SigO[C]$ quantifies the information
lost by choosing \emph{this} clock.  Different clocks
(different QRFs) give different $\SigO$---precisely
the result of De~Vuyst et~al.~\cite{DeVuyst2025}.


%=======================================================================
\section{Five Limits}
\label{sec:limits}
%=======================================================================

\subsection{Limit~1: Quantum information (Paper~I)}
\label{sec:limit_QI}

\textsc{Status: proved.}

\emph{Boundary conditions.}---Flat spacetime
($g_{\mu\nu} = \eta_{\mu\nu}$), finite-dimensional Hilbert
space $\hilb_S \otimes \hilb_E$, general CPTP map
$\chan: \mathcal{B}(\hilb_S) \to \mathcal{B}(\hilb_S)$.

\emph{Reduction.}---With no gravitational degrees of freedom,
$\rhost = |0\rangle\langle 0|$ (Minkowski vacuum) and
$\Sigma_{\mathrm{grav}} = 0$.  The only remaining
contribution is the channel-induced entropy production:
\begin{equation}
\Sigma_{\mathrm{QI}} \;=\;
\KL(\rho\|\sigma) - \KL\!\bigl(\chan(\rho)\|\chan(\sigma)\bigr)\,.
\label{eq:sigma_QI}
\end{equation}

\emph{What survives.}---The Petz recovery map
$\Petz_{\sigma,\chan}$, the JRSWW bound
$F \geq e^{-\Sigma/2}$~\cite{JRSWW2018,Fawzi2015},
and the equivalence chain of Paper~I:
\begin{equation}
\tauasym = 0
\;\Leftrightarrow\; \Sigma = 0
\;\Leftrightarrow\; \text{QMC}
\;\Leftrightarrow\; \text{KL conditions.}
\label{eq:equiv_chain}
\end{equation}

\emph{Mathematical precision.}---Exact.  No approximations.
Paper~I is a rigorous special case of
Eq.~\eqref{eq:master}~\cite{Paper1}.


\subsection{Limit~2: Strong-field gravity (Paper~II)}
\label{sec:limit_gravity}

\textsc{Status: derived (gravitational thermal attenuator
channel; saturation at $O((\rs/r)^2)$).}

\emph{Boundary conditions.}---Static, spherically symmetric,
asymptotically flat spacetime with mass $M$ at the origin.
Observer at infinity.

\emph{Reduction.}---Three independent
routes~\cite{Paper2}---modular flow
(Dorau--Much extension), gravitational Landauer
(Herrera~\cite{Herrera2020}), and pure-loss bosonic
channel---all yield the same result:
\begin{equation}
\Sgrav \;=\; -\ln(-g_{00})\,.
\label{eq:Sgrav}
\end{equation}
For the exponential metric
$g_{00} = -e^{-\rs/r}$~\cite{Yilmaz1958}:
$\Sgrav = \rs/r$ exactly.  For Schwarzschild:
$\Sgrav = -\ln(1 - \rs/r) = \rs/r + O((\rs/r)^2)$.

\emph{The saturation ansatz.}---If the JRSWW bound is
exactly saturated, $F = e^{-\Sigma/2}$, then
$\sqrt{-g_{00}} = e^{-\Sgrav/2}$, giving
$g_{00} = -e^{-\rs/r}$.  This ansatz is supported by
the triple identification
$\ceff/c = e^{-\Sgrav/2} = \Fbound$,
but exact saturation for $\Sigma > 0$ has no proof.
The approximate saturation gap is
$O\!\bigl((\rs/r)^2\bigr)$ in the weak
field~\cite{Paper2}.

\emph{Predictions.}---No event horizons
($\tauasym < 1$ everywhere); GW echoes at
$\Delta t \approx 4\,\rs/c$; shadow 4.6\% larger than
Schwarzschild; ISCO 5.6\% larger; QNM frequencies shifted
$\sim$4.4\%.

\emph{Open problem.}---No canonical gravitational CPTP
map $\chan_{\mathrm{grav}}$ has been constructed from
first principles whose QRE drop equals
$-\ln(-g_{00})$~\cite{Paper2}.


\subsection{Limit~3: Galactic dynamics (Paper~III)}
\label{sec:limit_galactic}

\textsc{Status: derived ($\eta = 1$ from DPI $+$ de~Sitter
extensivity; CMB remains open).}

\emph{Boundary conditions.}---Weak field, quasi-static,
$r \sim \mathrm{kpc}$.

\emph{Reduction.}---At galactic scales, the
renormalization-group (RG) flow of Newton's constant
introduces a logarithmic correction.  With marginal
anomalous dimension $\eta = 1$~\cite{Kumar2025}:
\begin{equation}
\Sigma(r) \;=\; \frac{\rs}{r}
\;+\; \frac{4\,k_*\,\rs}{\pi}\,
\ln\!\biggl(\frac{r}{r_0}\biggr)\,,
\label{eq:sigma_galactic}
\end{equation}
where $k_*$ is the crossover momentum scale and
$r_0$ a reference radius.

The DPI applied to the RG flow (Casini--Huerta
theorem~\cite{CasiniHuerta2017}) guarantees
\begin{equation}
\KL(\rho_{\mathrm{UV}}\|\sigma_{\mathrm{UV}})
\;\geq\;
\KL(\rho_{\mathrm{IR}}\|\sigma_{\mathrm{IR}})\,,
\label{eq:DPI_RG}
\end{equation}
so coarse-graining reduces the QRE---the running of $G$
is a DPI applied to the master equation.

The logarithmic correction produces a constant circular
velocity $v^2 = 2\GN M k_*/\pi$, reproducing flat
rotation curves.  Gubitosi et~al.~\cite{Gubitosi2024}
demonstrate competitive fits to 100 SPARC
galaxies~\cite{SPARC}.

\emph{CMB resolution via Khronon condensate
[structural framework established; full Boltzmann code
verification pending].}---Running $G$ alone cannot
explain CMB acoustic peaks: without cold dark matter at
the background level, $z_{\mathrm{eq}} \sim 530$ instead
of 3400.  However, the OEE postulate
(Sec.~\ref{sec:OEE}) promotes $u^a$ to a dynamical
Khronon field.  The Khronon mass
$\mu = 2\pi(1+\Omega_{\mathrm{b}})/r_{\mathrm{d}}$
(Sec.~\ref{sec:running_mu})---determined by
recombination physics, not by $H_0$---provides a
dust-like component at recombination that restores
the correct $z_{\mathrm{eq}}$ without particle dark
matter.  The Khronon connects to
AeST/Khronon-type
theories~\cite{SkordisZlosnik2021,BlanchetSkordis2024}
and provides CDM-like behavior at the perturbative
level via the ghost condensation condition $c_s^2 = 0$
(see Sec.~\ref{sec:running_mu} for details), as
confirmed by the BS2025 CMB fit~\cite{BS2025}.


\subsection{Limit~4: Cosmological expansion (FRW)}
\label{sec:limit_cosmo}

\textsc{Status: partially proved.}

\emph{Boundary conditions.}---Homogeneous, isotropic,
FRW spacetime.  The Kodama vector~\cite{Kodama1980}
$K^a = a^{-1}(1,\,-Hr,\,0,\,0)$ replaces the Killing
vector.

\emph{Background.}---The Clausius relation at the apparent
horizon $\rA = 1/H$, combined with the
Bekenstein--Hawking entropy $S_A = \pi/(\GN H^2)$ and
temperature $T_A = H/(2\pi)$, yields the Friedmann
equations~\cite{CaiKim2005,Jacobson1995}:
\begin{align}
H^2 &= \frac{8\pi\GN}{3}\,\rho\,,
\label{eq:friedmann1}\\
\dot{H} &= -4\pi\GN(\rho + p)\,.
\label{eq:friedmann2}
\end{align}
In the language of Eq.~\eqref{eq:master}, the background
satisfies $\Sigma_{\mathrm{background}} = 0$:
entanglement equilibrium~\cite{Jacobson2016}.

\emph{First-order perturbation theory.}---Using the
Aalsma--Bak modular Hamiltonian for FRW causal
diamonds~\cite{AalsmaBak2025},
$\delta K = 2\SdS\,\Phi$, together with the entanglement
first law $\delta D^{(1)} = \delta\langle K\rangle
- \delta S$~\cite{Faulkner2014}, one obtains
\begin{equation}
\delta D^{(1)} \;=\; S_A\,(2\Phi + \delta_m) \;=\; 0\,,
\label{eq:first_order_D}
\end{equation}
which reproduces the sub-Hubble Poisson equation
$\Phi = -\tfrac{1}{2}\delta_m$.

\emph{The modular Khronon.}---The perturbation of
modular time away from coordinate time defines
\begin{equation}
\delta\Tmod \;=\; -\frac{2\Phi}{H}\,,
\label{eq:modular_khronon}
\end{equation}
a scalar field encoding the deviation of the observer's
modular flow from the Kodama flow.  In the matter era,
$\Phi = \mathrm{const}$ (growing mode), so $\delta\Tmod$
is non-propagating: $\omega = 0$, $c_s = 0$.
This is the same dispersion relation as the
Blanchet--Skordis Khronon~\cite{BlanchetSkordis2024}.

\emph{Dark energy interpretation.}---Padmanabhan's CosMIn
principle~\cite{Padmanabhan2017} requires that the
total cosmic information remains finite:
$\int_0^\infty \tauasym(t)\,dt < \infty$.  Without a
cosmological constant, the integrand grows as $t$
increases (matter domination), and the integral diverges.
Accelerated expansion ($\Lambda > 0$) dilutes the matter
contribution, ensuring convergence.  In the
$\tauasym$ language: the universe \emph{must} eventually
stop producing entropy at an unbounded rate, and
$\Lambda > 0$ achieves this.

\emph{Ghost condensation resolves the epoch
problem.}---The modular Khronon inherits its behavior
from $\Phi$ and therefore has $c_s \neq 0$ in the
radiation era.  The independent Khronon kinetic term
$K(Q) = \mu^2(Q-1)^2$ ensures $c_s = 0$ at all epochs
via the ghost condensation condition $K'(Q_0) = 0$
(Paper~III).
The Khronon mass
$\mu = 2\pi(1+\Omega_{\mathrm{b}})/r_{\mathrm{d}}$
(Sec.~\ref{sec:running_mu}) is fixed at recombination
and provides CDM-like behavior at all subsequent
epochs.  The $c_s = 0$ condition ensures dust-like
perturbation growth independent of the background
equation of state.


\subsection{Limit~5: Observer-dependent time (Paper~V)}
\label{sec:limit_observer}

\textsc{Status: proved (core theorems).}

\emph{Boundary conditions.}---General system, arbitrary
observer $\mathcal{O}$ with algebra
$\Aobs \subseteq \Atot$.

\emph{Reduction.}---The observer $\mathcal{O}$ defines an
effective channel
$\chan_{\mathcal{O}}: \Atot \to \Aobs$
(conditional expectation), giving
\begin{equation}
\SigO \;=\; \KL(\omega\|\omega_{\mathrm{ref}})
\;-\; \KL\!\bigl(\chan_{\mathcal{O}}(\omega)\|
\chan_{\mathcal{O}}(\omega_{\mathrm{ref}})\bigr)\,.
\label{eq:sigma_O}
\end{equation}

Five results~[\textsc{proved} except where noted]:

\begin{theorem}[Monotonicity]
\label{thm:monotone}
If $\mathcal{O}_1 \leq \mathcal{O}_2$ (i.e.,
$\Aobs[1] \subseteq \Aobs[2]$), then
$\tauO[1] \geq \tauO[2]$.
More access implies a weaker arrow of time.
\end{theorem}

\begin{theorem}[Threshold]
\label{thm:threshold}
For each $\mathcal{O}$, there exists
$\epsilon_{\mathcal{O}} > 0$ such that
$\tauO < \epsilon_{\mathcal{O}}$ implies the process is
operationally time-reversible for $\mathcal{O}$.
\end{theorem}

\begin{theorem}[Subadditivity]
\label{thm:subadditivity}
$\tauasym_{\mathcal{O}_1 \vee \mathcal{O}_2}
\leq \min(\tauO[1],\,\tauO[2])$.
Combining observers always helps.
\end{theorem}

\begin{theorem}[Complementary uncertainty]
\label{thm:complementary}
For a complete set of $d+1$ MUB observers and any pure state:
$\sum_i \tauO[i] \geq (d+1) - \sqrt{2(d+1)} > 0$.
One cannot eliminate the arrow of time in all
``directions'' simultaneously.
(See Paper~V for proof via the 2-design purity sum rule.)
\end{theorem}

\begin{theorem}[Observer-dependent $\Sigma$]
\label{thm:obs_sigma}
$\SigO[1] \geq \SigO[2]$
whenever $\mathcal{O}_1 \leq \mathcal{O}_2$.
\end{theorem}

\emph{Connection to the master equation.}---The observer
$\mathcal{O}$ selects which subalgebra of $\Atot$ to
restrict to.  Different choices correspond to different
physical communities:
\begin{center}
\begin{tabular}{ll}
Lab physicist & $\to$ Paper~I \\
Distant astronomer & $\to$ Paper~II \\
Galactic surveyor & $\to$ Paper~III \\
Cosmic observer & $\to$ Paper~IV (cosmological) \\
Any finite agent & $\to$ Paper~V \\
Total observer & $\to$ $\tauasym = 0$ (Wheeler--DeWitt)
\end{tabular}
\end{center}

\emph{Fundamental barriers to $\tauasym \to 0$.}---Three
structural barriers prevent any observer from achieving
$\tauasym = 0$:
(i)~quantum uncertainty ($\tauasym_{\mathrm{quantum}}
\sim e^{-2\Delta D}$, Heisenberg);
(ii)~chaos ($\tauasym_{\mathrm{chaos}} \sim
e^{\lambda_L t}\epsilon$, Lyapunov);
(iii)~computational irreducibility
($\tauasym_{\mathrm{irreducible}} > 0$, Turing).
These provide structural limits on retrodiction,
connecting quantum mechanics, chaos theory, and
computability through a single parameter.  We
emphasize that this connection is conceptual,
not quantitative.


\subsection{The lattice structure unifying all limits}
\label{sec:lattice_unification}

The five limits are not independent theories patched
together but a single equation restricted to different
subalgebras.  The observers form a lattice under
refinement, and the DPI guarantees that moving up the
lattice (gaining access to more degrees of freedom)
always decreases $\tauO$.

Table~\ref{tab:limits} summarizes the five limits
and makes explicit the boundary conditions, the
mathematical reduction, and the physical
output.

\begin{table*}[t]
\caption{\label{tab:limits}The five limits of
the master equation~\eqref{eq:master}.  Each
row corresponds to a choice of observer algebra
$\Aobs$ and boundary conditions.  The third column
gives the resulting entropy production, and the
fourth column lists the central physical output.}
\begin{ruledtabular}
\begin{tabular}{llllc}
\textbf{Limit} &
\textbf{Boundary conditions} &
\textbf{$\Sigma$ formula} &
\textbf{Physical output} &
\textbf{Status} \\
\colrule
QI (Paper~I) &
Flat spacetime, finite $\hilb$, CPTP $\chan$ &
$\KL(\rho\|\sigma) - \KL(\chan\rho\|\chan\sigma)$ &
Petz recovery, QEC equivalence chain &
\textsc{proved} \\
Gravity (Paper~II) &
Static, $r \gg \rs$ &
$-\ln(-g_{00})$ &
Exponential metric, no horizon &
\textsc{derived} \\
Galactic (Paper~III) &
Weak field, $r \sim$ kpc &
$\rs/r + (4k_*\rs/\pi)\ln(r/r_0)$ &
Flat rotation curves &
\textsc{derived} \\
Cosmological &
FRW, Kodama vector &
$\delta D^{(1)} = S_A(2\Phi + \delta_m)$ &
Friedmann eqs., Khronon ($\mu = 2\pi(1\!+\!\Omega_{\mathrm{b}})/r_{\mathrm{d}}$) &
\textsc{derived} \\
Observer (Paper~V) &
General, $\Aobs \subseteq \Atot$ &
$\KL(\omega\|\omega_{\mathrm{ref}}) -
\KL(\chan_{\!\mathcal{O}}\omega\|
\chan_{\!\mathcal{O}}\omega_{\mathrm{ref}})$ &
Monotonicity, threshold, subadditivity &
\textsc{proved} \\
\end{tabular}
\end{ruledtabular}
\end{table*}

The lattice structure implies a nesting
of effective theories.  The thermodynamic
observer $\mathcal{O}_{\mathrm{thermo}}$
(access to coarse classical variables) sits
at the bottom; adding phase-space resolution
(momenta, quantum coherences, gravitational
wave channels) progressively refines the
observer toward $\top$.  Each refinement
step respects the DPI:
$\tauO[\mathrm{coarse}]
\geq \tauO[\mathrm{fine}]$.
Conversely, each coarse-graining step
creates an apparent ``cost of ignorance''
that manifests as entropy production.

This structure provides a concrete
realization of the complementarity between
communities.  The QI and GR communities are
not studying different physics---they are
studying the same $\Sigma$ restricted to
different subalgebras.  The proposed framework
provides a concrete formulation of this observation.


%=======================================================================
\section{Observer Entanglement Equilibrium}
\label{sec:OEE}
%=======================================================================

\subsection{The postulate}

The observer's 4-velocity $u^a$ enters the master equation
through the choice of foliation: the surfaces orthogonal
to $u^a$ define the observer's ``now.''  In the standard
QRE formalism (Sec.~\ref{sec:limit_cosmo}), $u^a$ is
fixed by the Kodama vector and is not varied.  We now
elevate it to a dynamical degree of freedom.

\begin{postulate}[Observer Entanglement Equilibrium (OEE)]
\label{post:OEE}
The physical state satisfies
\begin{align}
\frac{\delta\Sigma}{\delta g_{ab}} &= 0
\quad\text{(Einstein equations)}\,,
\label{eq:OEE_metric}\\[4pt]
\frac{\delta\Sigma}{\delta u^a} &= 0
\quad\text{(observer field equation)}\,,
\label{eq:OEE_observer}
\end{align}
subject to the unit-norm constraint $u^a u_a = -1$.
\end{postulate}

Equation~\eqref{eq:OEE_metric} reproduces gravity
(Dorau--Much~\cite{DorauMuch2025}).
Equation~\eqref{eq:OEE_observer} is \emph{our new
postulate}: it states that the QRE must be stationary
with respect to variations of the observer's reference
frame.  This postulate does not follow from known
results and represents the principal new hypothesis of
this work.

\subsection{Physical content}

The OEE postulate replaces an external choice
(``which observer?'') with a dynamical equation
(``the observer is self-consistently determined'').
This has three consequences.

\emph{Consequence~1: $u^a$ becomes the aether.}---A
unit timelike vector field $u^a$ with its own field
equation is precisely the Einstein-aether
field~\cite{Jacobson2004}.  In the hypersurface-orthogonal
limit $u_a = -N\,\partial_a T$, it is the Khronon
field~\cite{Jacobson2010,BlanchetSkordis2024}.
The identification
\begin{equation}
u^a\;\text{(observer)} \;\equiv\; A^a\;\text{(aether)}
\;\equiv\; n^a\;\text{(Khronon normal)}
\label{eq:identification}
\end{equation}
is exact at the level of mathematical objects: they are
all unit timelike vector fields.

\emph{Consequence~2: Petz optimality selects $K(Q)$.}---The
Petz recovery map is the \emph{unique} optimal retrodiction
(Paper~I).  Applied to perturbations of the observer field,
Petz optimality requires that background perturbations
minimize information loss---i.e., the kinetic function
$K(Q)$ must have a \emph{quadratic minimum} at the
background value $Q_0 = 1$:
\begin{equation}
K(Q) \;=\; \mu^2\,(Q - 1)^2 + O\!\bigl((Q-1)^3\bigr)\,.
\label{eq:K_Q}
\end{equation}
This is the Blanchet--Skordis kinetic
term~\cite{BlanchetSkordis2024}.  The argument is
heuristic, not rigorous: Petz optimality constrains
the \emph{form} of $K(Q)$ but not the mass scale $\mu$.

\emph{Consequence~3: $c_s = 0$.}---A kinetic function with
a quadratic minimum at $Q_0$ produces a scalar perturbation
with vanishing sound speed~\cite{BlanchetSkordis2024}:
\begin{equation}
\omega^2 \;=\; 0
\quad\text{(non-propagating mode).}
\label{eq:cs_zero}
\end{equation}
The perturbation grows under gravitational instability
without oscillating---exactly like pressureless cold dark
matter.  This $c_s = 0$ condition, combined with the
GW170817 constraint $c_T = c$ (requiring
$c_1 + c_3 = 0$~\cite{GW170817}), reduces the
Einstein-aether theory to a one-parameter family
after imposing $c_1 + c_2 + c_3 = 0$.

\subsection{Connection to AeST}

The Skordis--Zlosnik AeST theory~\cite{SkordisZlosnik2021}
is the only relativistic MOND theory that simultaneously
satisfies $c_T = c$, fits the CMB power spectra, and
reproduces MOND at galactic scales.  Its field content---a
unit timelike vector $A^\mu$, a shift-symmetric scalar
$\phi$, and the metric $g_{\mu\nu}$---maps onto the
OEE framework as follows:
\begin{center}
\begin{tabular}{lll}
\textbf{OEE} & \textbf{AeST} & \textbf{Match} \\
\hline
$u^\mu$ & $A^\mu$ & Exact \\
$\Sgrav$ & $\phi_s$ & Proportional in weak field \\
$K(Q)$ & $\mathcal{F}(Y,Q)$ & Same form near minimum \\
$\mu = 2\pi(1\!+\!\Omega_{\mathrm{b}})/r_{\mathrm{d}}$ [ours] & $\mu$ (free) & OEE proposes it \\
OEE [our postulate] & (imposed) & OEE provides motivation \\
\end{tabular}
\end{center}

The OEE postulate provides a physical interpretation---the
aether is the observer's time direction---and a principle
that selects the form of the action.  Within this
framework, the Khronon mass scale
$\mu = 2\pi(1+\Omega_{\mathrm{b}})/r_{\mathrm{d}}$
is determined by recombination physics
(Sec.~\ref{sec:running_mu}), matching the BS-fitted
value to $0.04\%$.  The free function
$\mathcal{F}(Y,Q)$ and the MOND interpolating function
$J(Y)$ are \emph{not} determined by the framework.

\subsection{Parameter constraints from observations and principles}

The Einstein-aether kinetic tensor
$K^{ab}_{cd}\nabla_a A^c \nabla_b A^d$ depends on four
dimensionless constants $c_1, c_2, c_3, c_4$.
Existing constraints reduce these:
\begin{enumerate}
\item GW170817~\cite{GW170817}: $|c_T/c - 1| < 3 \times
  10^{-15}$ requires $c_1 + c_3 = 0$.
\item CMB ($c_s = 0$): $c_1 + c_2 + c_3 = 0$, combined
  with the above, gives $c_2 = 0$.
\item Petz optimality: $K(Q)$ has a quadratic minimum,
  selecting the \emph{form} of the action but not
  numerical values.
\item DPI: $c_1 + c_4 \geq 0$ (positive energy
  condition), a sign constraint only.
\end{enumerate}
After imposing all constraints:
$c_2 = 0$, $c_1 = -c_3$ (free), and $c_4 > -c_1$.
In the Khronon limit (hypersurface-orthogonal $u^a$),
the mass scale $\mu$ is now determined:
$\mu = 2\pi(1+\Omega_{\mathrm{b}})/r_{\mathrm{d}}$
(Sec.~\ref{sec:running_mu}), fixed by recombination
physics.  One coupling constant
($c_4$, subject to $c_4 > -c_1$) remains
undetermined.

\subsection{What OEE does not determine}
\label{sec:OEE_limitations}

We must be explicit about the limitations:
\begin{enumerate}
\item $\OmDM h^2 \approx 0.12$ is an
  \emph{order-of-magnitude prediction}
  (Sec.~\ref{sec:running_mu}).  The Khronon mass
  $\mu = 2\pi(1+\Omega_{\mathrm{b}})/r_{\mathrm{d}}$
  determines the mass scale but
  $\OmDM \sim \rho_{\mathrm{crit}}/3 \approx 0.33$
  remains approximate; the exact $O(1)$ prefactor
  requires the full nonlinear AeST calculation.
\item The MOND scale $a_0 \approx 1.2 \times 10^{-10}$\,m/s$^2$
  is \emph{predicted to $0.3\%$}
  (Sec.~\ref{sec:running_mu}).
  Equation~\eqref{eq:a0_pred} gives
  $a_0 = 1.197 \times 10^{-10}\,\mathrm{m\,s}^{-2}$
  from CMB-measured quantities alone.  The sub-percent
  correction $(1+\Omega_{\mathrm{b}})$ has not been
  derived from first principles.
\item The coupling constant $c_4$ remains a free parameter.
  The Khronon mass scale $\mu$ is now determined
  ($\mu = 2\pi(1+\Omega_{\mathrm{b}})/r_{\mathrm{d}}$;
  Sec.~\ref{sec:running_mu}), but $c_4$ is constrained
  only by $c_1 + c_4 \geq 0$ (positive energy).
\item The free function $\mathcal{F}(Y,Q)$ of AeST
  determines the MOND interpolating function $J(Y)$
  and is chosen to fit observations.  While the DPI
  suggests that $J(Y)$ should follow from the RG
  flow of the gravitational channel (via the
  Casini--Huerta monotonicity
  theorem~\cite{CasiniHuerta2017}), no explicit
  derivation exists.
\item The transition from cosmological (CDM-like) to
  galactic (MOND-like) behavior is encoded in the
  kinetic function $K(Q)$ and its DBI
  completion~\cite{BlanchetSkordis2024}, not derived
  from the QRE.
\end{enumerate}

\subsection{Embedding vs.\ derivation}

To be precise about the status: the OEE postulate
[our new postulate], combined with the Khronon mass
determination of Sec.~\ref{sec:running_mu} [our result],
provides a physically motivated \emph{partial embedding}
of AeST/Khronon physics within the $\tauasym$ framework.
The Khronon mass
$\mu = 2\pi(1+\Omega_{\mathrm{b}})/r_{\mathrm{d}}$
is an empirical relation matching the BS-fitted value
to $0.04\%$; a first-principles derivation from the
Khronon action is not yet available.
The $O(1)$ prefactor in $\OmDM$ and the full nonlinear
AeST action are not determined.
The partial derivation would become \emph{complete} if
any of the following could be shown:
(i)~$D(\rhost\|\rhom)$ perturbation theory on FRW
naturally produces a Khronon-like field with the
correct kinetic term;
(ii)~the $O(1)$ prefactor in
$\OmDM \sim \rho_{\mathrm{crit}}/3$ is fixed by the
normalization of the Khronon kinetic function;
(iii)~the ghost condensation mechanism locks onto the
sound horizon at recombination, deriving
Relation~II from the Khronon action;
(iv)~the constrained optimization of
$D(\rhost\|\rhom)$ produces the Khronon-Tensor field
content~\cite{BlanchetSkordis2025} as Lagrange
multipliers.
Items~(i)--(iv) remain open calculations.


\subsection{The Khronon mass from recombination}
\label{sec:running_mu}

The kinetic function
$K(Q) = \mu^2(Q-1)^2$~[Eq.~\eqref{eq:K_Q}] contains
the mass scale $\mu$ as a free parameter.
In an earlier version of this paper, we proposed
$\mu_0 = H_0/c$ from dimensional analysis of the
de~Sitter thermal scale.
This identification is now \textbf{excluded}:
CLASS/Boltzmann calculations show that $\mu = H_0/c$
(whether constant or running as $\mu(a) = H(a)/c$)
is incompatible with CMB data at $>55\sigma$, because
it forces $\delta_0 \sim O(1)$ and hence
$\tilde{w}_0 \sim 0.1$---far too large.
The Blanchet--Skordis (BS) fitted value
$\mu^{-1} = 22.3\,\mathrm{Mpc}$~\cite{BS2025}
passes the CMB; the naive $c/H_0 \approx 4450\,
\mathrm{Mpc}$ overshoots by a factor of $\sim$200.

We now report two numerical relations
[\textsc{new synthesis}; first presented in
Ref.~\cite{Huang2026_mu}]
that pin down $\mu$ in terms of independently measured
cosmological quantities with sub-percent accuracy.

\emph{Relation~I: MOND acceleration at decoupling.}---We
propose
\begin{equation}
\mu\, c^2 \;=\; a_0\,(1 + z_{\mathrm{dec}}) \,,
\label{eq:mu_rel1}
\end{equation}
where $z_{\mathrm{dec}} = 1089.92$ is the photon
decoupling redshift~\cite{Planck2018}.
Inverting gives
$\mu^{-1} = c^2/[a_0(1+z_{\mathrm{dec}})]
= 22.25\,\mathrm{Mpc}$, matching the BS fit to
$0.2\%$.
Physically, the Khronon mass equals the MOND
acceleration blueshifted to the epoch of last
scattering: $\mu$ is fixed at the cosmological
moment when baryonic perturbations become
dynamically independent.

\emph{Relation~II: Baryon-loaded sound horizon.}---%
Independently,
\begin{equation}
\mu \;=\; \frac{2\pi\,(1 + \Omega_{\mathrm{b}})}
{r_{\mathrm{d}}} \,,
\label{eq:mu_rel2}
\end{equation}
where $r_{\mathrm{d}} = 147.09\,\mathrm{Mpc}$ is the
comoving sound horizon at the baryon drag
epoch~\cite{Planck2018} and
$\Omega_{\mathrm{b}} = 0.0493$.
Inverting:
$\mu^{-1} = r_{\mathrm{d}}/[2\pi(1+\Omega_{\mathrm{b}})]
= 22.31\,\mathrm{Mpc}$, matching the BS fit to
$0.04\%$.
The Khronon Compton wavelength is the reduced acoustic
wavelength $r_{\mathrm{d}}/(2\pi)$, corrected by the
baryon loading factor $(1+\Omega_{\mathrm{b}})$.
The Khronon field develops its condensate mass when
the baryon--photon plasma decouples; the fundamental
Fourier mode $k_{\mathrm{fund}} = 2\pi/r_{\mathrm{d}}$
sets the minimal mass for CDM-like clustering across
the full causal volume.

\emph{Consistency and the MOND prediction.}---The two
relations are structurally independent (Relation~I
uses $a_0$ and $z_{\mathrm{dec}}$; Relation~II
uses $r_{\mathrm{d}}$ and $\Omega_{\mathrm{b}}$)
yet agree to $0.27\%$.
Eliminating $\mu$ yields a parameter-free prediction
for the MOND acceleration:
\begin{equation}
\boxed{\;
a_0 \;=\; \frac{2\pi\,(1+\Omega_{\mathrm{b}})\,c^2}
{r_{\mathrm{d}}\,(1+z_{\mathrm{dec}})} \;}
\label{eq:a0_pred}
\end{equation}
Evaluating with Planck~2018 values:
\begin{equation}
a_0^{\mathrm{pred}}
= \frac{2\pi \times 1.0493 \times
        (2.998 \times 10^8)^2}
       {4.539 \times 10^{24} \times 1090.92}
= 1.197 \times 10^{-10}\;\mathrm{m\,s}^{-2},
\label{eq:a0_numerical}
\end{equation}
matching the empirical MOND value
$a_0 = (1.20 \pm 0.02_{\mathrm{stat}}
\pm 0.24_{\mathrm{sys}}) \times 10^{-10}\,
\mathrm{m\,s}^{-2}$~\cite{McGaugh2016,Lelli2017}
to $0.3\%$---a $30$--$40\times$ improvement over
Milgrom's $a_0 = cH_0/(2\pi)$ ($13\%$ off) and
Verlinde's $a_0 = cH_0/6$ ($9\%$ off).

\emph{Why $\mu_0 = H_0/c$ fails.}---The old
identification gave $a_0 = cH_0/(2\pi)
= 1.04 \times 10^{-10}\,\mathrm{m\,s}^{-2}$
(8\% off), and suffered from two fatal problems:
(i)~it has no theoretical justification from the
$\Sigma$ framework beyond dimensional analysis;
(ii)~CLASS calculations confirm that $\mu = H_0/c$
(constant or running) forces the Khronon equation of
state $\tilde{w}_0 \sim 0.13$, which is excluded by
CMB data.
Relation~\eqref{eq:mu_rel2} explains the factor-of-200
hierarchy: $\mu/(H_0/c) = (1+z_{\mathrm{dec}})/(2\pi)
\times a_0/a_0^{\mathrm{M}} \approx 200$, where
$a_0^{\mathrm{M}} \equiv cH_0/(2\pi)$.
The physics of recombination inserts a factor
$(1+z_{\mathrm{dec}})/(2\pi)$ that cannot be
neglected.

\emph{The $(1+\Omega_{\mathrm{b}})$ puzzle.}---At
leading order, $\mu = 2\pi/r_{\mathrm{d}}$ matches
the BS value to $5\%$.  The correction
$(1+\Omega_{\mathrm{b}})$ improves this to $0.05\%$.
Three plausible mechanisms exist at
$O(\Omega_{\mathrm{b}})$:
(a)~baryon gravitational self-energy correction to
the condensate mass;
(b)~the effective sound horizon without baryons
exceeds $r_{\mathrm{d}}$ by a factor
${\sim}(1+\Omega_{\mathrm{b}})$, so
$\mu = 2\pi/r_{\mathrm{d}}^{(0)}$ with
$r_{\mathrm{d}}^{(0)} \approx
r_{\mathrm{d}}(1+\Omega_{\mathrm{b}})$;
(c)~numerical coincidence at the $0.1\%$ level.
A definitive discrimination requires computing
the Khronon mass from the full perturbation theory
on the FRW background including baryonic effects.
We classify this sub-percent correction as
\textbf{unproven}.

\emph{Connection to $\Sigma$.}---On FRW backgrounds
with the Khronon condensate,
$\Sigma = 2\ln Q$ where $Q = 1 + \delta$.
The ghost condensation condition $K'(Q_0) = 0$ ensures
$c_s^2 = 0$ exactly (Paper~III).
Equation~\eqref{eq:a0_pred} adds a new element:
$a_0$ is the acceleration below which the
retrodiction channel (Petz recovery) transitions from
reversible to irreversible, with the transition scale
set at the epoch of baryon--photon decoupling.

\emph{The gravitational seesaw.}---An intriguing
numerical coincidence links $\mu$ to the neutrino
mass scale:
\begin{equation}
\mu \;\approx\; \frac{(\sum m_\nu)^2}{M_{\mathrm{Pl}}}\,,
\label{eq:seesaw}
\end{equation}
where $\sum m_\nu \approx 0.06\,\mathrm{eV}$ is
the minimal neutrino mass sum and
$M_{\mathrm{Pl}}$ is the reduced Planck mass.
The agreement is $1.7\%$.
The ghost condensation scale
$M = \sqrt{\mu\, M_{\mathrm{Pl}}} \approx 59\,
\mathrm{meV}$
is tantalizingly close to the lightest neutrino mass.
Whether this is coincidence or a genuine seesaw-type
mechanism remains open.

\emph{The unification chain (revised).}---Recombination
physics determines everything:
\begin{equation}
r_{\mathrm{d}},\,\Omega_{\mathrm{b}},\,z_{\mathrm{dec}}
\;\longrightarrow\; \mu = \frac{2\pi(1+\Omega_{\mathrm{b}})}{r_{\mathrm{d}}}
\;\longrightarrow\;
\begin{cases}
a_0 = \dfrac{\mu\,c^2}{1+z_{\mathrm{dec}}}
  & \text{(MOND scale, 0.3\%)}\\[4pt]
\OmDM \sim \rho_{\mathrm{crit}}/3
  & \text{(DM fraction)}
\end{cases}
\label{eq:unification_chain}
\end{equation}
The MOND scale is no longer $a_0 \sim cH_0$ but
$a_0 = 2\pi(1+\Omega_{\mathrm{b}})c^2/
[r_{\mathrm{d}}(1+z_{\mathrm{dec}})]$---an
\emph{imprint of recombination}, not a reflection of
the present-day Hubble rate.
This formula is immune to the $H_0$ tension because
the product $r_{\mathrm{d}}(1+z_{\mathrm{dec}})$
depends on the physical densities
$\Omega_{\mathrm{b}} h^2$ and
$\Omega_{\mathrm{m}} h^2$, not on $h$ alone.
The dark matter fraction $\OmDM \sim 1/3$ remains
an order-of-magnitude estimate ($\sim$25\% above
observed); the exact $O(1)$ prefactor depends on
the normalization of $K(Q)$.

\emph{Implications for the BS theory.}---If
Relation~II holds as an identity rather than a
coincidence, it reduces the free parameter count
of the BS Khronon theory from two ($a_0$, $\mu$) to
one ($a_0$ alone, with $\mu$ derived).
If the consistency condition~\eqref{eq:a0_pred} is
exact, $a_0$ itself is determined by cosmological
parameters, leaving the theory with \emph{zero} free
parameters in the dark sector beyond $\Lambda$CDM
inputs.

\emph{Honest assessment.}---Relations~\eqref{eq:mu_rel1}
and \eqref{eq:mu_rel2} are empirical findings with
sub-percent accuracy, not first-principles derivations.
The leading-order match ($\mu \approx 2\pi/r_{\mathrm{d}}$,
$5\%$) is physically motivated by the Fourier structure
of the sound horizon.
The sub-percent correction $(1+\Omega_{\mathrm{b}})$
and the underlying reason for the sound-horizon matching
both require a microscopic derivation from the Khronon
action.
The probability of accidental agreement at the $0.3\%$
level is $\sim 1/300$ for a single relation; accounting
for the independent leading-order match and the correct
order of the correction, the joint probability of
coincidence is $\sim 1/30{,}000$---strong evidence for
a physical connection, though not a proof.


%=======================================================================
\section{The Modular Khronon}
\label{sec:khronon}
%=======================================================================

\subsection{Derivation from FRW QRE perturbation theory}

We summarize the calculation detailed in the companion
research notes~\cite{Paper4Notes}.

On the FRW background, the comoving observer has
$u^a = (1,0,0,0)$, and $\Sigma_{\mathrm{grav}} =
-\ln(-g_{00}) = -\ln(1) = 0$.  At the apparent horizon
$\rA = 1/H$, the Kodama vector vanishes and
$\tauasym \to 1$.

At first order in perturbation theory, the modular
Hamiltonian of the FRW causal diamond~\cite{AalsmaBak2025}
receives a correction
\begin{equation}
\delta K \;=\; 2\SdS\,\Phi\,,
\label{eq:delta_K}
\end{equation}
where $\SdS = \pi/(\GN H^2)$.  The entanglement first law
$\delta D^{(1)} = \delta\langle K\rangle - \delta S_A = 0$
reproduces the linearized Poisson equation.

The modular Khronon $\delta\Tmod$ is defined as the
deviation of modular time from coordinate time.  From the
KMS condition applied to the perturbed state:
\begin{equation}
\delta\Tmod \;=\; -\frac{\delta K}{\langle K\rangle}
\;\cdot\; \frac{1}{\kappa}
\;=\; -\frac{2\Phi}{H}\,,
\label{eq:delta_Tmod}
\end{equation}
where $\kappa = H$ is the surface gravity of the FRW
apparent horizon.

\subsection{Dispersion relation}

In the matter-dominated era, the growing mode has
$\Phi = \mathrm{const}$, so $\delta\Tmod$ is frozen:
\begin{equation}
\omega \;=\; 0\,,\qquad c_s^2 \;=\; 0
\quad\text{(matter era).}
\label{eq:dispersion_matter}
\end{equation}
The modular Khronon perturbation is effectively
pressureless---it grows in place without oscillating,
exactly like CDM density perturbations.

In the radiation-dominated era, $\Phi$ oscillates and
decays inside the horizon, so $\delta\Tmod$ inherits
oscillatory behavior with an effective
$c_s^2 \neq 0$.  This is a departure from the
Blanchet--Skordis Khronon, which has $c_s = 0$ at
\emph{all} epochs due to the independent kinetic term
$K(Q)$.

\subsection{What works and what does not}

\begin{center}
\begin{tabular}{lccc}
\textbf{Feature} &
\textbf{Modular Khronon} &
\textbf{BS Khronon} &
\textbf{Running $\mu$} \\
\hline
$c_s = 0$ (matter era) & Yes & Yes & Yes \\
$c_s = 0$ (radiation era) & No & Yes & Yes \\
Background energy & 0 & Tuned & $\sim\!\rho_{\mathrm{crit}}/3$ \\
Coupling to $\Phi$ & Inherited & Independent & Derived \\
Mass parameter & $\sim M_{\mathrm{Pl}}$ & $\sim H_0$ (free) & $2\pi(1\!+\!\Omega_{\mathrm{b}})/r_{\mathrm{d}}$ (proposed) \\
\end{tabular}
\end{center}

The recombination-based determination
(Sec.~\ref{sec:running_mu}) offers a resolution of
the hierarchy problem:
$\mu = 2\pi(1+\Omega_{\mathrm{b}})/r_{\mathrm{d}}$
is set by the sound horizon, not put in by hand.
The factor of $\sim$200 between $\mu$ and $H_0/c$ is
explained by $(1+z_{\mathrm{dec}})/(2\pi)$---a
geometric-cosmological ratio determined by
recombination physics.

\subsection{Second-order effects and the positivity constraint}

At second order in perturbation theory, the QRE receives a
contribution
\begin{equation}
\delta D^{(2)} \;=\; \frac{1}{2}\,S_A\,
\bigl(4\Phi^2 + \delta_m^2\bigr)
\;+\; \delta^{(2)}\!\langle K\rangle
\;-\; \delta^{(2)} S_A\,.
\label{eq:second_order}
\end{equation}
The DPI requires $\delta D^{(2)} \geq 0$, which places a
\emph{sign constraint} on the second-order entanglement
entropy perturbation.  In the language of the modular
Khronon, this translates to the requirement that the
perturbation $\delta\Tmod$ does not grow faster than the
gravitational potential allows:
\begin{equation}
\abs{\delta\Tmod}^2 \;\leq\; \frac{S_A}{2}\,
\bigl(2\Phi + \delta_m\bigr)^2\,.
\label{eq:positivity}
\end{equation}
This is the information-theoretic analog of the positive
energy condition: the observer's time perturbation is
bounded by the available gravitational entropy.

The physical content is important.  In standard
cosmological perturbation theory, the growing mode
$\delta_m \propto a(t)$ is unconstrained by any
information-theoretic bound.  In the $\tauasym$ framework,
the DPI provides a universal upper bound on
structure growth---a prediction that could in principle
be tested against N-body simulations, though the bound is
far from saturated in the observed universe.

\subsection{The stealth black hole tension}

AeST admits ``stealth'' black hole solutions with
\emph{exact} Schwarzschild
geometry~\cite{SkordisVokrouhlicky2024}: the aether and
scalar carry non-trivial secondary hair while the metric
is identical to GR.  By contrast, Paper~II predicts the
exponential metric $g_{00} = -e^{-\rs/r}$ (no event
horizon).

These are observationally distinguishable.
If the OEE postulate embeds the $\tauasym$ framework
into AeST, then AeST's stealth solutions would supersede
Paper~II's exponential metric in the strong field.  Three
resolutions are possible:
(i)~the exponential metric is the correct weak-field
approximation while the stealth BH is the strong-field
completion;
(ii)~the $\tauasym$ framework should embed in a modified
AeST where stealth solutions do not exist;
(iii)~the identification breaks down in the strong-field
regime and quantum gravity effects dominate.
This tension is real and must be resolved by
experiment---the echo predictions of Paper~II versus the
Schwarzschild QNM spectrum of standard AeST.

\subsection{Connection to Blanchet--Skordis (2024) Khronon theory}

The Blanchet--Skordis
action~\cite{BlanchetSkordis2024,BlanchetSkordis2025}
\begin{equation}
S \;=\; \frac{c^3}{16\pi G}\int\!\sqrt{-g}\,
\bigl[R - 2J(Y) + 2K(Q)\bigr]\,d^4x
\;+\; S_{\mathrm{matter}}
\label{eq:BS_action}
\end{equation}
uses the Khronon scalar $\tau$ (whose gradient gives
$n^\mu$), acceleration-squared $Y$, and the lapse-like
variable $Q = c\sqrt{-g^{\mu\nu}\partial_\mu\tau\,
\partial_\nu\tau}$.

The OEE postulate [our new postulate] provides a
physical motivation for this action: $J(Y)$ encodes
the MOND interpolating function (from the DPI structure
of the gravitational channel at galactic scales), and
$K(Q)$ has a quadratic minimum (from Petz optimality
of the background retrodiction).  The mass scale $\mu$
is determined [our result]:
$\mu = 2\pi(1+\Omega_{\mathrm{b}})/r_{\mathrm{d}}$
(Sec.~\ref{sec:running_mu}), matching the BS-fitted
value to $0.04\%$; a first-principles derivation
from the Khronon action is pending.  The functional
form of $J(Y)$ is not determined.


%=======================================================================
\section{Falsifiable Predictions}
\label{sec:predictions}
%=======================================================================

The following predictions require the combined structure
of Papers~I--V and cannot be made by any individual paper
alone.  We list them as numbered, falsifiable statements.

\begin{enumerate}

\item \textbf{Joint echo + rotation curve test.}---
Paper~II predicts GW echoes at
$\Delta t \approx 4\,\rs/c$.  Paper~III predicts flat
rotation curves from running $G$ with a universal
crossover scale $k_* \sim 2.7 \times 10^{-2}\,
\mathrm{kpc}^{-1}$.  The unified framework requires
both to hold \emph{simultaneously}---the same
$\Sgrav = -\ln(-g_{00})$ that gives the echo prediction
must, at galactic scales with RG corrections, give flat
rotation curves.  \emph{Falsification:} detection of
echoes without the predicted rotation curve behavior
(or vice versa) would falsify the unified framework
while leaving the individual papers intact.
\emph{Timeline:} LIGO O5+ data (echoes); ongoing
SPARC analysis (rotation curves).

\item \textbf{$\tauasym$ consistency across scales.}---
For any physical system at any scale,
\begin{equation}
\tauO \;=\; 1 - F\!\bigl(\rho,\;
\Petz_{\mathcal{O}}\!\circ\!\chan_{\mathcal{O}}(\rho)\bigr)
\label{eq:consistency}
\end{equation}
must satisfy: (i)~$\tauO \in [0,1]$;
(ii)~monotonicity under refinement;
(iii)~$\tauasym = 0$ iff quantum Markov;
(iv)~$\tauasym = 1 - \sqrt{-g_{00}}$ for gravitational
observers; (v)~inclusion of RG-running corrections at
galactic scales; (vi)~approach to 1 at the apparent
horizon.
\emph{Falsification:} failure of any condition at any
scale falsifies the unified framework.

\item \textbf{Observer-dependent apparent dark matter.}---
Theorem~\ref{thm:obs_sigma} implies that different
observers assign different $\SigO$ to the same galaxy.
The ``missing mass'' attributed to dark matter is
\begin{equation}
\Delta\Sigma \;=\; \SigO[{\mathrm{limited}}]
\;-\; \SigO[{\mathrm{full}}] \;\geq\; 0\,.
\label{eq:apparent_DM}
\end{equation}
Quantum sensors with access to gravitational entanglement
degrees of freedom should, in principle, see less apparent
dark matter than classical telescopes.
\emph{Falsification:} quantum-gravitational measurements
showing no reduction in inferred dark matter.
\emph{Status:} qualitative and currently far from
testable, but structurally unique to the framework.

\item \textbf{Gravitational bound on retrodiction
fidelity.}---The framework implies that retrodiction
fidelity for any observer is bounded below by the
gravitational entropy production at that observer's
location:
\begin{equation}
\tauO \;\geq\; \tauasym_{\mathrm{grav}}
\;=\; 1 - e^{-\Sgrav/2}\,.
\label{eq:retrodiction_bound}
\end{equation}
\emph{Falsification:} laboratory demonstration of
retrodiction fidelity exceeding the gravitational
bound.

\item \textbf{MOND interpolating function from
Crooks.}---The Crooks fluctuation theorem
$P(\Sigma)/P(-\Sigma) = e^{\Sigma}$~\cite{Crooks1999}
suggests the interpolating function
\begin{equation}
\mu(x) \;=\; 1 - e^{-x}\,,\qquad
x \;\equiv\; a/a_0\,,
\label{eq:mu_crooks}
\end{equation}
which gives $\mu \approx x$ for $x \ll 1$ (deep MOND)
and $\mu \approx 1$ for $x \gg 1$ (Newtonian).
This form transitions faster than the standard
$\mu(x) = x/(1+x)$.
\emph{Falsification:} SPARC rotation curve
fits~\cite{SPARC} rejecting this functional form.
\emph{Status: speculative.}  A proper derivation
connecting $\Sigma$ to the effective gravitational
coupling is needed.

\item \textbf{MOND scale from recombination.}---
The sound horizon formula (Sec.~\ref{sec:running_mu})
gives
\begin{equation}
a_0 \;=\; \frac{2\pi(1+\Omega_{\mathrm{b}})\,c^2}
{r_{\mathrm{d}}\,(1+z_{\mathrm{dec}})}
\;\approx\; 1.197 \times 10^{-10}\;\mathrm{m/s}^2\,,
\label{eq:a0_KMS}
\end{equation}
matching the empirical MOND value
$1.2 \times 10^{-10}\;\mathrm{m/s}^2$ to $0.3\%$.
This supersedes the earlier KMS-based estimate
$a_0 = cH_0/(2\pi) \approx 1.04 \times 10^{-10}\,
\mathrm{m/s}^2$ ($13\%$ off) and Verlinde's
$a_0 = cH_0/6$ ($9\%$ off)~\cite{Verlinde2016}.
\emph{Falsification:} precise $a_0$ measurements
deviating from $1.20 \times 10^{-10}\;\mathrm{m/s}^2$
by more than the ${\sim}\,20\%$ systematic uncertainty.
\emph{Status:} empirical with $0.3\%$ accuracy,
not a derivation.

\end{enumerate}


%=======================================================================
\section{Honest Assessment}
\label{sec:honest}
%=======================================================================

Table~\ref{tab:honest} summarizes the epistemic status of
every major claim.  The framework's genuine contribution is
\emph{synthesis}: connecting known results across communities.
Approximately 60\% of the content is established results,
25\% is new synthesis, 10\% is new formulas or definitions,
and 5\% is new interpretation.

\begin{table*}[t]
\caption{\label{tab:honest}Epistemic status of claims in
the unified framework.}
\begin{ruledtabular}
\begin{tabular}{llll}
\textbf{Claim} &
\textbf{Status} &
\textbf{What would change it} &
\textbf{Paper} \\
\colrule
$\tauasym = 0 \Leftrightarrow \Sigma = 0
\Leftrightarrow$ QMC $\Leftrightarrow$ QEC
& \textsc{proved} & Mathematical counterexample & I \\
$F \geq e^{-\Sigma/2}$ (JRSWW)
& \textsc{proved} & --- (theorem) & I \\
$\Sgrav = -\ln(-g_{00})$ (three routes)
& \textsc{derived} & Alternative channel models & II \\
Exponential metric from saturation
& \textsc{conjectured} & Prove/disprove saturation for $\Sigma > 0$ & II \\
GW echoes at $4\,\rs/c$
& \textsc{predicted} & LIGO O5+ data & II \\
Flat curves from running $G$
& \textsc{derived} & DPI $+$ extensivity $\Rightarrow$ $\eta=1$ & III \\
Friedmann from $\delta D = 0$
& \textsc{proved} & --- (Jacobson/Cai--Kim) & IV \\
$\delta\Tmod = -2\Phi/H$
& \textsc{derived} & Independent verification & IV \\
OEE postulate [\textsc{our postulate}]
& \textsc{new} & Derive from first principles & IV \\
$u^a \equiv A^a$ (observer $=$ aether)
& \textsc{structural} & --- (same math object) & IV \\
$K(Q)$ quadratic from Petz
& \textsc{heuristic} & Rigorous proof & IV \\
$\mu = 2\pi(1\!+\!\Omega_{\mathrm{b}})/r_{\mathrm{d}}$ [\textsc{our result}]
& \textsc{empirical, $0.04\%$} & Derive from Khronon action & IV \\
$\mu c^2 = a_0(1+z_{\mathrm{dec}})$ [\textsc{our result}]
& \textsc{empirical, $0.2\%$} & Derive from first principles & IV \\
$\OmDM \sim \rho_{\mathrm{crit}}/3$ [\textsc{our estimate}]
& \textsc{order-of-mag.} & Fix $O(1)$ prefactor from nonlinear AeST & IV \\
$a_0 = 2\pi(1\!+\!\Omega_{\mathrm{b}})c^2/[r_{\mathrm{d}}(1\!+\!z_{\mathrm{dec}})]$ [\textsc{our result}]
& \textsc{$0.3\%$ match} & Derive $(1\!+\!\Omega_{\mathrm{b}})$ factor & IV \\
$\mu_0 = H_0/c$ [\textsc{excluded}]
& \textsc{dead ($>55\sigma$)} & --- & IV \\
$\tauO$ monotonicity
& \textsc{proved} & --- (DPI) & V \\
Fundamental barriers to $\tauasym \to 0$
& \textsc{structural} & Quantitative bounds on retrodiction & V \\
\end{tabular}
\end{ruledtabular}
\end{table*}


%=======================================================================
\section{What This Framework Is Not}
\label{sec:not}
%=======================================================================

Honest delimitation is essential.

\emph{Not a competitor to string theory or LQG.}---Those
are microscopic theories seeking a fundamental Lagrangian.
The $\tauasym$ framework is an organizational principle,
analogous to thermodynamics.  If string theory is
statistical mechanics, the $\tauasym$ framework is
thermodynamics: it constrains without specifying
microscopic dynamics.

\emph{Not a microscopic theory.}---It does not specify
the fundamental degrees of freedom, the UV completion,
or the Planck-scale physics.  It uses the information-theoretic
properties of QRE (positivity, DPI, monotonicity) that
hold for \emph{any} microscopic theory satisfying
unitarity.

\emph{Does not precisely predict $\OmDM h^2$.}---The
Khronon mass
$\mu = 2\pi(1+\Omega_{\mathrm{b}})/r_{\mathrm{d}}$
determines the mass scale but yields only an
order-of-magnitude estimate
$\OmDM \sim \rho_{\mathrm{crit}}/3 \approx 0.33$,
which is $\sim$25\% above the observed value of 0.265.
The discrepancy lies in the undetermined $O(1)$
prefactor of the kinetic function $K(Q)$.  The
\emph{order of magnitude} is a genuine prediction;
the precise value requires the full nonlinear AeST
calculation, which this framework does not perform.

\emph{Does not prove the exponential metric.}---The
saturation of $F = e^{-\Sigma/2}$ is an ansatz.
The Thattarampilly--Zheng--Kakkat
solution~\cite{Thattarampilly2026} of the full GfE
equations gives Schwarzschild + $O(r^{-4})$ corrections,
\emph{not} the exponential metric.  Whether the
correct realization of the master equation produces the
exponential metric or Schwarzschild + corrections
remains open.

\emph{Does not replace the full CMB calculation.}---The
Khronon mass
$\mu = 2\pi(1+\Omega_{\mathrm{b}})/r_{\mathrm{d}}$
matches the BS2025 value that passes the CMB
(Sec.~\ref{sec:running_mu}), but the full CMB acoustic
peak structure requires running the perturbation
equations through a Boltzmann code with the Khronon
source terms---a calculation that has been done for
AeST~\cite{SkordisZlosnik2021} and by
BS2025~\cite{BS2025}, but not yet with the $\mu$
value derived from our Relations~I--II.  The pathway
is clear; the detailed computation remains to be
performed.


%=======================================================================
\section{Open Problems}
\label{sec:open}
%=======================================================================

We list the open problems in decreasing order of
severity.

\begin{enumerate}
\item \emph{The canonical gravitational CPTP map.}
  No channel $\chan_{\mathrm{grav}}$ has been
  constructed from first principles with QRE drop
  equal to $-\ln(-g_{00})$.  This is the most serious
  gap in Paper~II and propagates to the unified
  framework.

\item \emph{Petz saturation for $\Sigma > 0$.}
  Exact saturation of the JRSWW bound is needed for
  the exponential metric.  No known channel saturates
  for $\Sigma > 0$.  The Golden--Thompson inequality
  provides a structural obstruction.

\item \emph{Remaining coupling constant $c_4$.}
  The Khronon mass scale $\mu$ is now determined
  ($\mu = 2\pi(1+\Omega_{\mathrm{b}})/r_{\mathrm{d}}$;
  Sec.~\ref{sec:running_mu}), but the Einstein-aether
  parameter $c_4$ remains free.  The DPI constrains
  $c_1 + c_4 \geq 0$ (positive energy) but nothing
  more.

\item \emph{$O(1)$ prefactor in $\OmDM$.}
  The framework predicts
  $\OmDM \sim \rho_{\mathrm{crit}}/3 \approx 0.33$,
  which is $\sim$25\% above the observed 0.265.  The
  discrepancy lies in the normalization of $K(Q)$,
  which requires the full nonlinear AeST calculation.

\item \emph{Microscopic derivation of
  $\mu = 2\pi(1+\Omega_{\mathrm{b}})/r_{\mathrm{d}}$.}
  The two numerical relations
  (Sec.~\ref{sec:running_mu}) match the BS-fitted
  value to sub-percent accuracy, but neither has been
  derived from the Khronon action.
  A microscopic derivation would need to show that
  ghost condensation locks onto the acoustic scale
  at recombination.  The $(1+\Omega_{\mathrm{b}})$
  correction also requires explanation.

\item \emph{The anomalous dimension $\eta_G = 1$.}
  Paper~III uses $\eta_G = 1$ (Kumar~\cite{Kumar2025})
  for flat rotation curves.  The DPI $+$ extensivity
  argument~\cite{Paper3} supports $\eta_G = 1$, but
  a fully rigorous proof from the gravitational RG
  remains open.

\item \emph{Wheeler--DeWitt as Level~0.}
  The statement $\tauasym_{\mathrm{universe}} = 0$
  (no global arrow of time) is motivated by the
  Wheeler--DeWitt equation $\hat{H}|\Psi\rangle = 0$,
  but inherits its well-known problems (operator
  ordering, inner product).  The statement
  $\tauasym_{\mathrm{total}} = 0$ may be
  independently justifiable without invoking WDW.

\item \emph{The precise bridge between Bianconi GfE
  and Araki QRE.}
  The GfE uses a generalized matrix relative entropy
  on metric tensors; the master equation uses the
  Araki relative entropy on von Neumann algebras.
  The exact mathematical relationship is unknown.
\end{enumerate}


%=======================================================================
\section{Discussion and Outlook}
\label{sec:discussion}
%=======================================================================

\subsection{The framework as thermodynamics}

The relationship between the $\tauasym$ framework and
a future microscopic theory of quantum gravity is
analogous to the relationship between thermodynamics
and statistical mechanics.  Thermodynamics provides
universal laws ($\Delta S \geq 0$,
$dU = TdS - PdV$) that constrain all systems without
specifying microscopic dynamics.  The $\tauasym$
framework provides universal laws
($\tauO \geq 0$, $F \geq e^{-\Sigma/2}$,
monotonicity under coarse-graining) that constrain
all quantum-gravitational systems.

The value of an organizational principle lies not in
its predictive power for specific numbers (the
framework will never predict the mass of the electron)
but in its ability to connect seemingly unrelated
phenomena.  That the same parameter $\tauasym$ governs
QEC decoder performance, gravitational redshift,
galactic rotation curves, and the cosmological arrow
of time is the content of this paper.

\subsection{Connections to existing programs}

The Petz recovery map sits at the intersection of
several programs.  We compare honestly, emphasizing
complementarity rather than competition.

\emph{String theory and LQG} are microscopic theories
seeking a fundamental Lagrangian.  The $\tauasym$
framework is a different \emph{type} of theory: it
constrains without specifying microscopic dynamics.
String theory provides UV completion; the $\tauasym$
framework provides universal organizational principles
that hold for any UV-complete theory satisfying
unitarity.  The relationship is analogous to
thermodynamics (universal laws) versus statistical
mechanics (microscopic derivation).

\emph{It from Qubit / holographic QEC:} The Petz
map reconstructs bulk operators from boundary data
in AdS/CFT~\cite{Cotler2019}.  Our framework
extends this beyond AdS, using relative entropy
(well-defined for Type~III algebras)
instead of entanglement entropy (which requires
holographic regularization).  The overlap is deep:
the same Petz map that reconstructs holographic
spacetime also quantifies the arrow of time.

\emph{ER = EPR~\cite{MaldacenaSusskind2013}:} A wormhole
connecting two entangled systems has $\tauasym = 0$ for
the total observer (the system is closed).  For a
one-sided observer, $\tauasym > 0$: the other side
appears thermal.  This is the ER = EPR
conjecture translated to $\tauasym$ language.
The recent operational proof by Kaplan et~al.\
(arXiv:2410.16496) that monogamous entanglement is
operationally indistinguishable from a wormhole is
consistent with our framework: the $\tauO$ values
assigned by separated observers to a shared entangled
state are consistent with a wormhole geometry.

\emph{Verlinde's emergent gravity~\cite{Verlinde2016}:}
Verlinde derives apparent dark matter from displaced
volume-law entanglement entropy.
Our running-$G$ mechanism (Paper~III) and the
recombination-based derivation of $a_0$
(Sec.~\ref{sec:running_mu})
are complementary pictures: Verlinde modifies the
force law, we modify the gravitational channel.  Both
agree that dark matter is an information-theoretic effect,
not new particles.  Our formula achieves $0.3\%$ accuracy
versus Verlinde's $9\%$.  Ghari and
Haghi~\cite{GhariHaghi2026} show Verlinde's prediction
is favored over MOND at $5.2\sigma$ for dwarf
spheroidals, providing observational support for the
general entropic-gravity program.

\emph{Padmanabhan's cosmic information
(CosMIn)~\cite{Padmanabhan2017}:}
The requirement that total cosmic information remains
finite, $I_{\mathrm{CosMIn}} = \int_0^\infty
\tauasym(t)\,dt < \infty$, requires $\Lambda > 0$:
accelerated expansion is needed to stop $\tauasym$
from growing indefinitely.  This provides a
thermodynamic argument for the cosmological
constant.

\emph{Wheeler's ``It from Bit''~\cite{Wheeler1990}:}
The framework offers one possible mathematical realization
of Wheeler's vision:
every physical observable is a function of
$\Sigma = \Daraki(\omega|_{\Aobs}\|
\omega_{\mathrm{ref}}|_{\Aobs})$.
The ``bit'' is the quantum relative entropy.
The ``it'' is emergent spacetime geometry, the arrow
of time, and the effective matter content.

\subsection{Central thesis}

The five papers, taken together, point toward a single
organizing statement:
\begin{quote}
\emph{Time is the cost of ignorance, quantified by
$\tauasym$.}
\end{quote}
The total observer ($\Aobs = \Atot$) has $\tauasym = 0$:
no arrow of time, no ignorance, no cost.
Every partial observer pays a price for not seeing
everything---this price is $\tauasym$, and it manifests
as thermodynamic irreversibility (Paper~I), gravitational
redshift (Paper~II), apparent dark matter (Paper~III),
cosmological expansion (this paper), and the subjective
flow of time (Paper~V).

Whether this is merely a useful organizing language or a
genuine insight into the structure of nature will be
decided by the predictions: GW echoes, rotation curves
from running $G$, the $\tauasym$ consistency condition
across scales.  The framework is falsifiable, and
that is all one can ask.


\subsection{Comparison with five existing programs}

To avoid giving the false impression that the
$\tauasym$ framework operates in a vacuum, we
compare concretely with five established programs,
emphasizing what each provides that the $\tauasym$
framework does not, and vice versa.

\emph{1.\ Jacobson's thermodynamic
gravity~\cite{Jacobson1995,Jacobson2016}.}---Jacobson
derives the Einstein equation from the Clausius
relation $\delta Q = T\,\delta S$ applied to local
Rindler horizons.  Our zeroth-order cosmological
limit (Sec.~\ref{sec:limit_cosmo}) is essentially
the FRW version of his argument, via Cai--Kim.
The $\tauasym$ framework adds: (i)~the JRSWW bound,
which gives a \emph{quantitative} fidelity guarantee
beyond the leading order; (ii)~the observer lattice
structure (Sec.~\ref{sec:lattice_unification});
(iii)~the OEE postulate that promotes the observer
to a dynamical field.  Jacobson's program does not
address dark matter phenomenology.

\emph{2.\ Verlinde's emergent
gravity~\cite{Verlinde2016}.}---Verlinde derives
apparent dark matter from displaced volume-law
entanglement entropy.  Our mechanism is complementary:
Paper~III uses running $G$ (a channel-level modification),
while Verlinde modifies the entropy budget.
Verlinde predicts $a_0 \sim cH_0/6$ ($9\%$ off);
our recombination-based formula
(Sec.~\ref{sec:running_mu}) achieves $0.3\%$.  Ghari and Haghi~\cite{GhariHaghi2026}
find that Verlinde's prediction is favored over
MOND at $5.2\sigma$ for dwarf spheroidals, which
is consistent with (but not specific to) our framework.
The $\tauasym$ framework adds: the QEC connection
(Paper~I), the exponential metric prediction
(Paper~II), and the observer-dependent formulation
(Paper~V).  Verlinde's program does not include these.

\emph{3.\ Skordis--Z\l{}o\'{s}nik
AeST~\cite{SkordisZlosnik2021}.}---AeST is the only
relativistic MOND theory that fits the CMB.  The
$\tauasym$ framework does not compete with AeST; it
provides a physical interpretation of AeST's field
content (the aether is the observer's time direction),
a selection principle for the kinetic function
(Petz optimality $\Rightarrow$ quadratic minimum),
and proposes the Khronon mass scale
$\mu = 2\pi(1+\Omega_{\mathrm{b}})/r_{\mathrm{d}}$
(Sec.~\ref{sec:running_mu}), matching the BS-fitted
value to $0.04\%$.
AeST provides: (i)~the full nonlinear action with
fitted parameters; (ii)~CMB power spectra;
(iii)~structure formation.  The $\tauasym$ framework
proposes a determination of the Khronon coupling that
AeST treats as a free parameter, and predicts the
MOND acceleration to $0.3\%$ from CMB quantities alone.

\emph{4.\ Bianconi's Gravity from
Entropy~\cite{Bianconi2025,BianconiCosmo2025}.}---The
GfE action is the most concrete realization of the
master equation.  Thattarampilly
et~al.~\cite{Thattarampilly2026} find that the
exact vacuum solution is Schwarzschild $+ O(r^{-4})$
corrections---\emph{not} the exponential metric.
This means our saturation ansatz (Paper~II) may be
inconsistent with the most natural GfE realization.
Whether an alternative GfE realization yields the
exponential metric is open.

\emph{5.\ It from Qubit / holographic
QEC~\cite{Cotler2019}.}---The Petz map reconstructs
bulk operators from boundary data in AdS/CFT.  Our
framework extends the Petz map beyond holographic
settings, using QRE (well-defined for Type~III
algebras) instead of entanglement entropy.  The
holographic program provides: UV completeness (in
the dual CFT), the RT formula, and explicit error
correction codes.  The $\tauasym$ framework provides:
the $\tauasym$ parameter connecting QEC to
thermodynamics, the Crooks--JRSWW bridge, and the
extension to cosmological settings where no AdS
boundary exists.

The overall picture is that the $\tauasym$
framework occupies a specific niche: it is more
general than any single program (connecting five
domains) but less predictive than any specific one.
This is the characteristic signature of an
organizational principle rather than a microscopic
theory.


%=======================================================================
\begin{acknowledgments}
%=======================================================================
The author thanks M.~M.~Wilde for correspondence on
quantum recovery bounds, and the anonymous readers whose
scrutiny improved the clarity of the honesty
classifications.
This work received no external funding.
\end{acknowledgments}


%=======================================================================
% REFERENCES
%=======================================================================
\begin{thebibliography}{99}

% Paper series
\bibitem{Paper1}
S.-K.~Huang,
``The Arrow of Time from Petz Recovery,''
(2026), Paper~I of this series.

\bibitem{Paper2}
S.-K.~Huang,
``Information Loss is Local: The Gravitational Refractive
Index as Petz Recovery Fidelity,''
(2026), Paper~II of this series.

\bibitem{Paper3}
S.-K.~Huang,
``Galactic Dynamics from Quantum Information Recovery,''
(2026), Paper~III of this series.

\bibitem{Paper5}
S.-K.~Huang,
``Observer-Dependent Temporal Asymmetry,''
(2026), Paper~V of this series.

\bibitem{Paper4Notes}
S.-K.~Huang,
``Perturbation Theory of $D(\rho_{\mathrm{spacetime}}
\|\rho_{\mathrm{matter}})$ on FRW Background,''
Research notes for Paper~IV (2026).

\bibitem{Huang2026_mu}
S.-K.~Huang,
``The Khronon Mass Scale from Recombination:
Two Numerical Relations Connecting MOND, CMB,
and Dark Matter,''
in preparation (2026).

% Petz map and recovery
\bibitem{Petz1988}
D.~Petz,
``Sufficiency of channels over von Neumann algebras,''
Q.~J.\ Math.\ \textbf{39}, 97 (1988).

\bibitem{Parzygnat2023}
A.~J.~Parzygnat and F.~Buscemi,
``Axioms for retrodiction: achieving time-reversal
symmetry with a prior,''
Quantum \textbf{7}, 1013 (2023).

\bibitem{JRSWW2018}
M.~Junge, R.~Renner, D.~Sutter, M.~M.~Wilde, and
A.~Winter,
``Universal recovery maps and approximate sufficiency
of quantum relative entropy,''
Ann.\ Henri Poincar\'{e} \textbf{19}, 2955 (2018).

\bibitem{Fawzi2015}
O.~Fawzi and R.~Renner,
``Quantum conditional mutual information and
approximate Markov chains,''
Commun.\ Math.\ Phys.\ \textbf{340}, 575 (2015).

% Crooks and fluctuation theorems
\bibitem{Crooks1999}
G.~E.~Crooks,
``Entropy production fluctuation theorem and the
nonequilibrium work relation for free energy
differences,''
Phys.\ Rev.\ E \textbf{60}, 2721 (1999).

% Araki and algebraic QFT
\bibitem{Araki1976}
H.~Araki,
``Relative entropy of states of von Neumann algebras,''
Publ.\ RIMS Kyoto \textbf{11}, 809 (1976).

\bibitem{Lindblad1975}
G.~Lindblad,
``Completely positive maps and entropy inequalities,''
Commun.\ Math.\ Phys.\ \textbf{40}, 147 (1975).

% Gravity from entropy / QRE
\bibitem{DorauMuch2025}
L.~Dorau and A.~Much,
``Deriving the Einstein equation from quantum
relative entropy,''
Phys.\ Rev.\ Lett.\ (2025);
arXiv:2510.24491.

\bibitem{Bianconi2025}
G.~Bianconi,
``Gravity from entropy,''
Phys.\ Rev.\ D \textbf{111}, 066001 (2025);
arXiv:2408.14391.

\bibitem{Jacobson1995}
T.~Jacobson,
``Thermodynamics of spacetime: The Einstein equation
of state,''
Phys.\ Rev.\ Lett.\ \textbf{75}, 1260 (1995).

\bibitem{Jacobson2016}
T.~Jacobson,
``Entanglement equilibrium and the Einstein equation,''
Phys.\ Rev.\ Lett.\ \textbf{116}, 201101 (2016).

\bibitem{Faulkner2014}
T.~Faulkner, M.~Guica, T.~Hartman, R.~C.~Myers,
and M.~Van~Raamsdonk,
``Gravitation from entanglement in holographic CFTs,''
J.\ High Energy Phys.\ \textbf{03}, 051 (2014).

% FRW thermodynamics
\bibitem{CaiKim2005}
R.-G.~Cai and S.~P.~Kim,
``First law of thermodynamics and Friedmann equations
of Friedmann--Robertson--Walker universe,''
J.\ High Energy Phys.\ \textbf{02}, 050 (2005).

\bibitem{Kodama1980}
H.~Kodama,
``Conserved energy flux for the spherically symmetric
system and the backreaction problem in the black hole
evaporation,''
Prog.\ Theor.\ Phys.\ \textbf{63}, 1217 (1980).

\bibitem{AalsmaBak2025}
L.~Aalsma and D.~Bak,
``Modular Hamiltonian for de~Sitter diamonds,''
Phys.\ Rev.\ D \textbf{112}, 026017 (2025);
arXiv:2503.04886.

% Einstein-aether and Khronon
\bibitem{Jacobson2004}
T.~Jacobson and D.~Mattingly,
``Einstein-aether theory,''
Phys.\ Rev.\ D \textbf{70}, 024003 (2004).

\bibitem{Jacobson2010}
T.~Jacobson,
``Extended Horava gravity and Einstein-aether theory,''
Phys.\ Rev.\ D \textbf{81}, 101502(R) (2010).

% AeST and Khronon theory
\bibitem{SkordisZlosnik2021}
C.~Skordis and T.~Z\l{}o\'{s}nik,
``New relativistic theory for Modified Newtonian
Dynamics,''
Phys.\ Rev.\ Lett.\ \textbf{127}, 161302 (2021).

\bibitem{BlanchetSkordis2024}
L.~Blanchet and C.~Skordis,
``A relativistic Khronon theory for Modified
Newtonian Dynamics,''
J.\ Cosmol.\ Astropart.\ Phys.\ \textbf{11}, 040 (2024);
arXiv:2404.06584.

\bibitem{BlanchetSkordis2025}
L.~Blanchet and C.~Skordis,
``Khronon-Tensor theory,''
(2025); arXiv:2507.00912.

% GW170817
\bibitem{GW170817}
B.~P.~Abbott \textit{et al.}\ (LIGO Scientific,
Virgo, Fermi~GBM, and INTEGRAL Collaborations),
``Gravitational waves and gamma-rays from a binary
neutron star merger: GW170817 and GRB~170817A,''
Astrophys.\ J.\ Lett.\ \textbf{848}, L13 (2017).

% Exponential metric
\bibitem{Yilmaz1958}
H.~Y\i lmaz,
``New approach to general relativity,''
Phys.\ Rev.\ \textbf{111}, 1417 (1958).

% Gravitational Landauer
\bibitem{Herrera2020}
L.~Herrera,
``Landauer principle and general relativity,''
Entropy \textbf{22}, 340 (2020).

% Running G and rotation curves
\bibitem{Kumar2025}
V.~Kumar,
``Rotation curves from quantum gravity,''
Phys.\ Lett.\ B \textbf{871}, 139503 (2025);
arXiv:2509.05246.

\bibitem{Gubitosi2024}
G.~Gubitosi \textit{et al.},
``Running Newton constant and rotation curves,''
(2024); arXiv:2403.00531.

\bibitem{SPARC}
F.~Lelli, S.~S.~McGaugh, and J.~M.~Schombert,
``SPARC: Mass models for 175 disk galaxies with
Spitzer photometry and accurate rotation curves,''
Astron.\ J.\ \textbf{152}, 157 (2016).

% Casini-Huerta RG = DPI
\bibitem{CasiniHuerta2017}
H.~Casini and M.~Huerta,
``Irreversibility of the renormalization group flow
in non-unitary quantum field theory,''
J.\ High Energy Phys.\ \textbf{06}, 002 (2017).

% Observer dependence
\bibitem{BassoCeleri2025}
M.~L.~W.~Basso, J.~Maziero, and L.~C.~Celeri,
``Quantum detailed fluctuation theorem in curved
spacetimes,''
Phys.\ Rev.\ Lett.\ \textbf{134}, 050406 (2025);
arXiv:2405.03902.

\bibitem{DeVuyst2025}
L.~De~Vuyst, S.~Eccles, P.~A.~H\"ohn, and
J.~Kirklin,
``Gravitational entropy is observer-dependent,''
J.\ High Energy Phys.\ (2025).

\bibitem{Chandrasekaran2023}
V.~Chandrasekaran, G.~Penington, and E.~Witten,
``Large $N$ algebras and generalized entropy,''
J.\ High Energy Phys.\ \textbf{04}, 009 (2023).

% Verlinde and MOND
\bibitem{Verlinde2016}
E.~P.~Verlinde,
``Emergent gravity and the dark universe,''
SciPost Phys.\ \textbf{2}, 016 (2017);
arXiv:1611.02269.

\bibitem{Milgrom2020}
M.~Milgrom,
``The $a_0$--cosmology connection in MOND,''
(2020); arXiv:2001.09729.

% Holographic QEC
\bibitem{Cotler2019}
J.~Cotler, P.~Hayden, G.~Penington, G.~Salton,
B.~Swingle, and M.~Walter,
``Entanglement wedge reconstruction via universal
recovery channels,''
Phys.\ Rev.\ X \textbf{9}, 031011 (2019).

% ER=EPR
\bibitem{MaldacenaSusskind2013}
J.~Maldacena and L.~Susskind,
``Cool horizons for entangled black holes,''
Fortschr.\ Phys.\ \textbf{61}, 781 (2013).

% Wheeler It from Bit
\bibitem{Wheeler1990}
J.~A.~Wheeler,
``Information, physics, quantum: The search for
links,'' in \emph{Complexity, Entropy, and the Physics
of Information}, edited by W.~H.~Zurek (Addison-Wesley,
1990).

% GfE solutions
\bibitem{Thattarampilly2026}
J.~Thattarampilly, X.~Zheng, and A.~Kakkat,
``Exact vacuum solutions in Gravity from Entropy,''
(2026); arXiv:2602.13694.

% Independent verification
\bibitem{Buscemi2024}
F.~Buscemi, J.~Bai, and V.~Scarani,
``Fully quantum stochastic entropy production,''
(2024); arXiv:2412.12489.

% Ghari-Haghi
\bibitem{GhariHaghi2026}
A.~Ghari and H.~Haghi,
``Verlinde's emergent gravity favored over MOND at
$5.2\sigma$ for dwarf spheroidals,''
(2026); arXiv:2601.01715.

% Padmanabhan CosMIn
\bibitem{Padmanabhan2017}
T.~Padmanabhan and H.~Padmanabhan,
``Cosmic information, the cosmological constant,
and the amplitude of primordial perturbations,''
Phys.\ Lett.\ B \textbf{773}, 81 (2017).

% Pazy entropic MOND
\bibitem{Pazy2013}
E.~Pazy,
``Quantum statistical modified entropic gravity as a
theoretical basis for MOND,''
Phys.\ Rev.\ D \textbf{87}, 084063 (2013).

% AeST analysis
\bibitem{Bataki2023}
Z.~Bataki, C.~Skordis, and T.~Z\l{}o\'{s}nik,
``Hamiltonian analysis of Aether Scalar Tensor theory,''
(2023); arXiv:2307.15126.

% Stealth BHs in AeST
\bibitem{SkordisVokrouhlicky2024}
C.~Skordis and D.~Vokrouhlick\'{y},
``Stealth black holes in Aether Scalar Tensor theory,''
(2024); arXiv:2412.15395.

% Page-Wootters
\bibitem{PageWootters1983}
D.~N.~Page and W.~K.~Wootters,
``Evolution without evolution: Dynamics described by
stationary observables,''
Phys.\ Rev.\ D \textbf{27}, 2885 (1983).

% McGaugh RAR
\bibitem{McGaugh2016}
S.~S.~McGaugh, F.~Lelli, and J.~M.~Schombert,
``Radial acceleration relation in rotationally supported
galaxies,''
Phys.\ Rev.\ Lett.\ \textbf{117}, 201101 (2016).

% Milgrom original
\bibitem{Milgrom1983}
M.~Milgrom,
``A modification of the Newtonian dynamics as a possible
alternative to the hidden mass hypothesis,''
Astrophys.\ J.\ \textbf{270}, 365 (1983).

% Smolin MOND from QG
\bibitem{Smolin2017}
L.~Smolin,
``MOND as a regime of quantum gravity,''
Phys.\ Rev.\ D \textbf{96}, 083523 (2017).

% Bianconi cosmological
\bibitem{BianconiCosmo2025}
G.~Bianconi,
``Gravity from Entropy: Cosmological solutions,''
(2025); arXiv:2510.22545.

% Kunz GDM
\bibitem{Kunz2009}
M.~Kunz,
``The dark degeneracy: On the number and nature of
dark components,''
Phys.\ Rev.\ D \textbf{80}, 123001 (2009).

% Pikovski
\bibitem{Pikovski2015}
I.~Pikovski, M.~Zych, F.~Costa, and \v{C}.~Brukner,
``Universal decoherence due to gravitational time
dilation,''
Nat.\ Phys.\ \textbf{11}, 668 (2015).

\bibitem{TKS2016}
S.~Thomas, M.~Kopp, and C.~Skordis,
``Constraining generalised dark matter parameters with
the CMB,''
Astrophys.\ J.\ \textbf{830}, 155 (2016);
arXiv:1601.05097.

\bibitem{Lelli2017}
F.~Lelli, S.~S.~McGaugh, J.~M.~Schombert, and
M.~S.~Pawlowski,
``One law to rule them all: The radial acceleration
relation of galaxies,''
Astrophys.\ J.\ \textbf{836}, 152 (2017).

\bibitem{Planck2018}
N.~Aghanim \textit{et al.} (Planck Collaboration),
``Planck 2018 results. VI. Cosmological parameters,''
Astron.\ Astrophys.\ \textbf{641}, A6 (2020);
arXiv:1807.06209.

\bibitem{BS2025}
L.~Blanchet and C.~Skordis,
``Khronon dark matter: CMB and large-scale structure,''
arXiv:2507.00912 (2025).

\end{thebibliography}


%=======================================================================
\appendix
\section{Derivation chain: FRW background to modular Khronon}
\label{app:derivation}
%=======================================================================

We collect the key steps of Sec.~\ref{sec:limit_cosmo}
and Sec.~\ref{sec:khronon} in a self-contained
derivation chain.

\emph{Step~1: Background state.}---On the flat FRW
background
$ds^2 = -dt^2 + a^2(t)\delta_{ij}\,dx^i\,dx^j$,
the apparent horizon sits at $\rA = 1/H$ with
associated thermodynamic quantities:
\begin{align}
T_A &= \frac{H}{2\pi}\,,\qquad
S_A = \frac{\pi}{\GN H^2}\,,\notag\\
E &= \frac{\rA}{2\GN} = \frac{1}{2\GN H}\,.
\label{eq:app_background}
\end{align}

\emph{Step~2: Thermal state identification.}---The
spacetime state at the apparent horizon is thermal
at the Kodama temperature:
$\rhost = Z_A^{-1}\exp(-K_A)$,
where the modular Hamiltonian is
$K_A = (2\pi/H)\cdot M(\rA) = S_A$
(Aalsma--Bak~\cite{AalsmaBak2025}).

\emph{Step~3: Zeroth-order QRE.}---At equilibrium
(Friedmann equations satisfied),
$D(\rhost\|\rhom) = 0$: the spacetime and matter
states are matched.  Out of equilibrium, the
Clausius relation $-dE = T_A\,dS_A$ projected along
the Kodama vector yields~\cite{CaiKim2005}
\begin{equation}
H^2 = \frac{8\pi\GN}{3}\,\rho\,.
\end{equation}

\emph{Step~4: First-order perturbation.}---In
Newtonian gauge
$g_{00} = -(1+2\Phi)$,
$g_{ij} = a^2(1-2\Psi)\delta_{ij}$
with $\Phi = \Psi$ (no anisotropic stress), the
modular Hamiltonian perturbation is
$\delta K = 2S_A\,\Phi$.
The entanglement first law gives
\begin{equation}
\delta D^{(1)} = \delta\langle K\rangle - \delta S_A
= S_A(2\Phi + \delta_m) = 0\,,
\end{equation}
reproducing $\Phi = -\delta_m/2$ (linearized
Poisson equation).

\emph{Step~5: Modular Khronon.}---Define the
deviation of modular time from coordinate time:
\begin{equation}
\delta\Tmod = -\frac{\delta K}{\langle K\rangle}
\cdot\frac{1}{\kappa}
= -\frac{2\Phi}{H}\,,
\end{equation}
where $\kappa = H$ is the surface gravity of
the FRW apparent horizon.

\emph{Step~6: Dispersion.}---In the matter era,
$\Phi = \mathrm{const}$ (growing mode), so
$\partial_t\delta\Tmod = 0$ and the mode has
$\omega = 0$, $c_s = 0$: a non-propagating
scalar perturbation that grows under gravitational
instability without oscillating.

\emph{Step~7: Identification.}---The field
$\delta\Tmod$ has the same algebraic structure and
dispersion relation as the Blanchet--Skordis
Khronon perturbation in the matter era.  The
physical identification is
$\delta\Tmod \equiv \delta T_{\mathrm{Khronon}}$,
with the caveat that the epoch-independent
$c_s = 0$ of the Blanchet--Skordis theory
requires an independent kinetic term $K(Q)$ that
the QRE formalism does not provide.


\end{document}
